% concatenated from Mikhovich_Subcontractible.tex
% user docs for compositio.cls
%\def\Filedate{2021/10/06} % date this file last revised
%\def\Fileversion{}    % version of compositio.cls documented

\documentclass[11pt,a4paper]{article}
\usepackage[T1]{fontenc}
\usepackage[utf8]{inputenc}
\usepackage[english]{babel}
\usepackage[margin=1in]{geometry}
\usepackage{amsmath}


%\usepackage{mathptmx}       % selects Times Roman as basic font
%\usepackage{helvet}         % selects Helvetica as sans-serif font
%\usepackage{courier}        % selects Courier as typewriter font
%\usepackage{type1cm}        % activate if the above 3 fonts are
                            % not available on your system
%
%\usepackage{makeidx}         % allows index generation
\usepackage{graphicx}        % standard LaTeX graphics tool
                             % when including figure files
%\usepackage{multicol}        % used for the two-column index
%\usepackage[bottom]{footmisc}% places footnotes at page bottom
%\usepackage[tbtags]{amsmath}

%\usepackage{mathtools}
\usepackage{amsfonts,amssymb}%,mathrsfs,amscd}
\usepackage[matrix,arrow,curve]{xy}
%\usepackage[xcolor,qtwo]{rvdtx}
%\usepackage{multicol}

%\usepackage{color}
%\usepackage{xspace}
%\usepackage{pdfwidgets}
%\usepackage{enumerate}
%\usepackage{mathptmx}
%\usepackage{amssymb}
\usepackage[numbers]{natbib}
\usepackage[hidelinks]{hyperref}
\usepackage{amsthm}

% If you have the AMSLaTeX distribution installed on your system,
% please delete the "[noams]" option above.

% definitions specific to this author guide only
\newcommand*{\code}[1]{{\mdseries\texttt{#1}}}
\newcommand*{\pkg}[1]{{\mdseries\textsf{#1}}}
\renewcommand{\topfraction}{0.99}
%\renewcommand{\contentsname}{Contents\\{\footnotesize\normalfont(A table
%of contents should normally not be included)}}
\newtheorem{theorem}{Theorem}[section] % number like 3.1, 3.2, 3.3, etc.
\newtheorem{definition}[theorem]{Definition} % numbered with thm
%\newtheorem{equaion}{Equation}
\newtheorem{lemma}[theorem]{Lemma}
\newtheorem{remark}[theorem]{Remark}
\newtheorem{proposition}[theorem]{Proposition}
\newtheorem{example}[theorem]{Example}
\newtheorem{corollary}[theorem]{Corollary}
%\newtheorem{proof}{Proof}

%
\begin{document}

\title{Bousfield--Kan Completions of Subcontractible Presentations}
%\
%\titlerunning{Rational and $p$-adic analogs of J.H.C. Whitehead's conjecture}
\author{Andrey Mikhovich\thanks{Moscow Center for Fundamental and Applied
Mathematics, Lomonosov Moscow State University, GSP-1, Leninskie Gory,
Moscow 119991, Russian Federation. Email: \texttt{amikhovich@gmail.com}.}}
\date{}
\maketitle

\begin{center}
\textit{Dedicated to Professor Alexander S. Mishchenko on the occasion of his
85th birthday, with deepest respect and admiration.}
\end{center}

%\institute{Andrey Mikhovich \at
 %             Moscow Center for Fundamental and Applied Mathematics, Moscow, Russia \\
      %        Tel.:+7 916 670 36 28\\
%              Fax.:+7 916 670 36 28\\
           %  \email{MikhovichAM@moaz.ru}           %  \\
 %            \emph{Present address:} of F. Author  %  if needed
 %        }
%\dedication{}

%\address{Moscow Center for Fundamental and Applied Mathematics}%
%Here\\MT55 9XX}
%\curraddr{Lomonosov Moscow State University, GSP-1, Leninskie Gory, Moscow, 119991, Russian Federation}
%\classification{55P60 (primary), 14F35, 18G40 (secondary). }

%\keywords{J.H.C. Whitehead's asphericity conjecture, cohomology of prounipotent groups, prounipotent crossed modules}
%Using these concepts we consider a prounipotent analog of J.H.C. Whitehead's asphericity conjecture and provide its positive solution.
%\date{Received: date / Accepted: date}
% The correct dates will be entered by the editor
%\thanks{This file documents \pkg{compositio} version \Fileversion\ and
%was last revised \Filedate.}



\begin{abstract}
We study Bousfield--Kan completions through the interaction of free
simplicial resolutions, their filtration spectral sequences, and a
noncommutative arithmetic square.  For every free discrete simplicial
group of finite type, we express its integral pronilpotent completion
as the homotopy pullback of its rational prounipotent completion and
the product of its pro-$p$ completions over an explicit adelic
simplicial group.  The adelic entry is formed by taking restricted
products at finite nilpotent stages and then their inverse limit;
no nilpotency assumption on the original group of components is required.
Finite subpresentations of contractible presentations provide an
explicit application of this construction.  Independence of the
specified relators makes the positive-degree terms of the rational
and mod-$p$ filtration spectral sequences vanish, with convergence
verified on the quotient towers.  Continuous comparison of free
simplicial resolutions then realizes, in characteristic zero, the
equivalence with a constant free prounipotent group by morphisms and
homotopies in that category.  For the corresponding presentation
complex $K$ we obtain $R_\infty K\simeq K(F_R(Z),1)$ for
$R=\mathbb Q,\mathbb F_p,\mathbb Z$, where $Z$ indexes a complementary
basis and $F_R(Z)$ denotes, respectively, the rational points of a
free prounipotent group, a free pro-$p$ group, or a free pronilpotent
group.  Compatible contractions at the nilpotent stages identify all
four entries of the arithmetic square in this case.
\end{abstract}
%\MSC{55P60 \sep 14F35 \sep 18G40}

% \PACS{PACS code1 \and PACS code2 \and more}

%\maketitle
%\newpageafter{abstract}
%\vspace*{6pt}\tableofcontents 

\section{Introduction}\label{s0}
The purpose of this paper is to describe how rational, pro-$p$, and
integral completion interact with the simplicial structure of a group
presentation.  Three constructions organize the argument: comparison
of free simplicial resolutions, the spectral sequences of their
lower-central and Zassenhaus filtrations, and a noncommutative
arithmetic square.  Finite subpresentations of contractible
presentations provide a concrete setting in which the interaction of
these constructions can be computed explicitly.  We retain the chosen
relators, construct comparison maps and homotopies in the completed
categories, and describe the passage from nilpotent quotients to the
integral completion.  The arithmetic-square construction itself
applies to every free discrete simplicial group of finite type.

A presentation $(X\mid Y)$ of a discrete group $G$ determines an exact sequence
\begin{equation}\label{eq1}
1 \rightarrow R \rightarrow \Phi \xrightarrow{\pi} G \rightarrow 1
\end{equation}
in which $\Phi=\Phi(X)$ is the free discrete group with a set $X$ of generators, and $R$
is a normal subgroup in $\Phi$ normally generated by a set $Y\subset R$ of defining relations. 
We use the notation $(X\mid Y)$ for a presentation with generating set $X$
and chosen set of relators $Y$.  The normal closure of $Y$ in $\Phi(X)$ will
be denoted by $R$.  Thus \eqref{eq1} is the extension determined by the
presentation, rather than the presentation itself.
By a subpresentation $(X_1\mid Y_1)$ we mean a choice of subsets
$X_1\subseteq X$ and $Y_1\subseteq Y$ such that the selected relators
are words in $\Phi(X_1)$.  A presentation is called \emph{aspherical}
when its associated two-dimensional CW-complex $K(X\mid Y)$ has
$\pi_2K(X\mid Y)=0$.  The construction is recalled in
Section~\ref{r1.2}; see also \cite{Hut}.

A presentation is called \emph{contractible} when its presentation complex is
contractible.  A \emph{subcontractible presentation} is a finite
subpresentation of a contractible presentation.

J.~H.~C. Whitehead's asphericity conjecture asserts that every connected
subcomplex of an aspherical two-dimensional CW-complex is aspherical; see
\cite{Ro} and the references therein.  Howie showed that it is enough to
consider subcomplexes of contractible two-complexes \cite{Ho83}.

\subsection{Presentations and their second homotopy groups}\label{r1.2}
Let $(X\mid Y)$ be a presentation of $G\cong\Phi(X)/R$, where $R$ is
the normal closure of the relators $r_y$, $y\in Y$.  Its standard
two-dimensional presentation complex $K(X\mid Y)$ has one $0$-cell, one
$1$-cell $e_x^1$ for each $x\in X$, and one $2$-cell $e_y^2$ attached along
a loop representing $r_y$ for each $y\in Y$.
Let $\widetilde K$ denote the universal cover.  Covering-space theory and the
Hurewicz theorem give natural isomorphisms
\[
 \pi_2K(X\mid Y)\cong\pi_2\widetilde K\cong H_2(\widetilde K;\mathbb Z).
\]

It turns out that the cellular chain complex $C_*(\widetilde K)$ is
$G$-isomorphic to the complex associated with the chosen presentation
\cite[Proposition~9]{BH1}.  We use right $\mathbb ZG$-modules throughout
this subsection and obtain the exact sequence
\begin{equation}\label{1.1}
0\longrightarrow\pi_2K(X\mid Y)\longrightarrow C_2(X\mid Y)
\xrightarrow{d_2}C_1(X\mid Y)\xrightarrow{d_1}C_0(X\mid Y)
\xrightarrow{\epsilon}\mathbb Z\longrightarrow0,
\end{equation}
where $C_0=\mathbb ZG$, $C_1=\bigoplus_{x\in X}e_x^1\mathbb ZG$,
and $C_2=\bigoplus_{y\in Y}e_y^2\mathbb ZG$.
Write $\pi$ also for the induced map $\mathbb Z\Phi\to\mathbb ZG$.
For $x\in X$, let $D_x^{\mathrm r}$ denote the right Fox derivative,
specified by
\[
 D_x^{\mathrm r}(x')=\delta_{x,x'},\qquad
 D_x^{\mathrm r}(uv)=D_x^{\mathrm r}(u)v+
                         \epsilon(u)D_x^{\mathrm r}(v)
 \quad(u,v\in\mathbb Z\Phi).
\]
The right-handed Fox identity is
\[
 a-\epsilon(a)=\sum_{x\in X}(x-1)D_x^{\mathrm r}(a)
 \qquad(a\in\mathbb Z\Phi);
\]
see \cite[Chapter~4]{BH1}, \cite[Section~4.2]{FENN}.
With these conventions the cellular boundaries are
\[
 d_1(e_x^1)=1-\pi(x),\qquad
 d_2(e_y^2)=\sum_{x\in X}e_x^1\,\pi(D_x^{\mathrm r}(r_y)).
\]
Indeed, $d_1d_2(e_y^2)=-\pi(r_y-1)=0$.
The image of $d_2$ is the relation module $R/[R,R]$, with right action
$\overline r\cdot g=\overline{\widetilde g^{-1}r\widetilde g}$ for
any lift $\widetilde g\in\Phi$ of $g\in G$.  Its Magnus embedding is
\[
 i:R/[R,R]\hookrightarrow\bigoplus_{x\in X}e_x^1\mathbb ZG,
 \qquad
 i(\overline r)=\sum_{x\in X}e_x^1\,\pi(D_x^{\mathrm r}(r));
\]
see \cite[Corollary~1]{BH1}.

For comparison, put $\Lambda=\mathbb Z\Phi$, let $\mathfrak f$ be
its augmentation ideal, and let
$\mathfrak r=\ker(\Lambda\to\mathbb ZG)$.
Gruenberg's free resolution begins
\begin{equation}\label{1.2}
\cdots\longrightarrow\mathfrak r^2/\mathfrak r^3
\longrightarrow\mathfrak f\mathfrak r/\mathfrak f\mathfrak r^2
\longrightarrow\mathfrak r/\mathfrak r^2
\longrightarrow\mathfrak f/\mathfrak f\mathfrak r
\longrightarrow\mathbb ZG\longrightarrow\mathbb Z\longrightarrow0;
\end{equation}
see \cite[Theorem~2]{Gru}.  To describe its degree-two free module,
choose a free group basis $B$ of the subgroup $R$.  Then
\[
 \mathfrak r=\bigoplus_{b\in B}(1-b)\Lambda,
 \qquad
 \mathfrak r^2=\bigoplus_{b\in B}(1-b)\mathfrak r,
\]
so the classes of $1-b$, $b\in B$, form a right $\mathbb ZG$-basis
of $\mathfrak r/\mathfrak r^2$.
Likewise, $\mathfrak f/\mathfrak f\mathfrak r$ is free on the classes
of $1-x$, $x\in X$.  The maps
\[
 \tau_0=\operatorname{id}_{\mathbb ZG},\qquad
 \tau_1(e_x^1)=1-x,\qquad \tau_2(e_y^2)=1-r_y
\]
give a comparison in degrees zero through two; commutation with the
differentials follows from the right Fox identity.  In general the
chosen relators $r_y$ are only normal generators of $R$ in $\Phi$,
and $\tau_2$ need not be an isomorphism.

For a contractible presentation, the required basis statement follows
directly from cellular chains.
\begin{lemma}\label{l.1.3}
Let $(X\mid Y)$ be a contractible presentation.  The images of the relators
$r_y$, $y\in Y$, form a basis of the free abelian group
$\Phi(X)_{\mathrm{ab}}\cong\mathbb Z^{(X)}$.  In particular, if $X$ and $Y$
are finite, then $|X|=|Y|$ and the exponent-sum matrix is unimodular.
\end{lemma}
\begin{proof}
Since $K(X\mid Y)$ is contractible and $G=1$, its augmented cellular chain
complex gives an exact sequence
\[
 0\longrightarrow\mathbb Z^{(Y)}\xrightarrow{\,\overline d_2\,}
 \mathbb Z^{(X)}\longrightarrow0.
\]
The matrix of $\overline d_2$ is the exponent-sum matrix
$\bigl(\epsilon(D_x^{\mathrm r}(r_y))\bigr)_{x,y}$.  Hence this homomorphism is an isomorphism.  In the finite case its matrix belongs to
$\mathrm{GL}_{|X|}(\mathbb Z)$.
\end{proof}
For a contractible presentation, $G=1$ and $R=\Phi$, so
$\mathfrak r/\mathfrak r^2\cong\Phi_{\mathrm{ab}}$.
The lemma therefore makes $\tau_2$ an isomorphism in this case,
without requiring the words $r_y$ to be a free group basis of $\Phi$.
%\begin{equation}\label{crossed}
%$$\xymatrix{\overline{C} \ar[r]^{\cong}\ar[d]^{\rho} & C_2(X|R) \ar[r]^{\cong} \ar[d]^{d_2} & C_2(\widetilde{K}) \ar[d]^{\partial_2} \\
%    \overline{R} \ar[r]^{i} & C_1(X|R) \ar[r]^{\cong} \ar[d]^{d_1} & C_1(\widetilde{K}) \ar[d]^{\partial_1}\\
%  & C_0(X|R) \ar[r]^{\cong} & C_0(\widetilde{K}) }$$
%\end{equation}
%This diagram shows, in particular, that $\pi=Ker\;\widetilde{\theta}\cong \pi_2(K(X\mid R)).$
\subsection{Simplicial presentation as a model of $\Omega K(X|R)$}\label{r1.1}
By definition, a simplicial resolution of a group $ G $ is a free simplicial group $F_{\bullet}$ (i.e. $F_n$ is a free group) such that $\pi_0F_{\bullet}\cong G$ and $\pi_nF_{\bullet}=0$ for $n>0$, this concept is analogous to the Eilenberg - MacLane space $ K (G, 1) $ in the category of simplicial groups. By \textbf{homotopy groups of a simplicial group} $F_{\bullet}$ we understand the homology groups of its \textbf{Moore complex} $(NF_n=\cap^{n-1}_{i=0}Ker\;d_i^n, d_n^n|_{NF_n})$.

The step-by-step construction of a free simplicial resolution starts
with generators in degree zero and adds generators in degree one for
the chosen relators; see \cite[Section~1.2]{MP}. Denote this initial
stage by $P_\bullet=F^{(1)}_\bullet$. It is a full simplicial group,
with no new nondegenerate free generators above degree one. Only
its restriction to degrees $0,1,2$ and its augmentation are displayed:
\begin{equation}\label{2}
\xymatrix@C=3.8em{
 P_2 \ar@<2ex>[r]^{d_0} \ar[r]|{d_1}
     \ar@<-2ex>[r]_{d_2} &
 P_1 \ar@<4ex>[l]^{s_0} \ar@<-4ex>[l]_{s_1}
     \ar@<1.5ex>[r]^{d_0} \ar@<-1.5ex>[r]_{d_1} &
 P_0 \ar@<3.5ex>[l]^{s_0} \ar[r] & G.}
\end{equation}
Here
\[
\begin{aligned}
 P_0&=\Phi(X),\qquad P_1=\Phi(s_0X\sqcup Y),\\
 P_2&=\Phi(s_1s_0X\sqcup s_0Y\sqcup s_1Y),
\end{aligned}
\]
where $s_0s_0=s_1s_0$, and the degree-one faces satisfy
\[
 d_0(s_0x)=d_1(s_0x)=x,\qquad d_0(y)=1,\qquad d_1(y)=r_y.
\]
Higher faces on the degenerate generators are determined by the
simplicial identities. In particular,
$\operatorname{im}(d_1:\ker d_0\to P_0)=R$ and $\pi_0P_\bullet=G$.
The diagram $P_1\rightrightarrows P_0$ with its degeneracy is called
a \textbf{simplicial presentation of $G$}.

The Moore complex of $P_\bullet$ has the form
\[
 \cdots\longrightarrow N_2P\xrightarrow{d_2}N_1P
 \xrightarrow{d_1}N_0P,
\]
where $N_0P=P_0$ and
$N_nP=\bigcap_{i=0}^{n-1}\ker(d_i:P_n\to P_{n-1})$ for $n\geq1$.
Thus
\[
 \pi_1P_\bullet
 =\frac{\ker(d_1:N_1P\to N_0P)}{d_2(N_2P)}.
\]
Generation by degeneracies in degree two does not imply $N_2P=1$.
In general $P_\bullet$ is not yet a resolution of $G$: its positive
homotopy groups have not been killed.

For the standard presentation complex $K=K(X\mid Y)$, choose Kan's
cellular free simplicial model $B_K$ with degree-zero generators $X$
and degree-one generators $Y$, with the attaching words specified
above \cite[Definition~5.2 and Theorem~6.2]{Kan}. There are no higher
nondegenerate generators because $K$ is two-dimensional. With these
choices $P_\bullet$ is this cellular model, and \eqref{2} displays its
restriction to degrees $0,1,2$. Kan's comparison
\cite[Theorem~5.5]{Kan} gives the loop homotopy equivalence with
$GS_1K$, where $S_1K$ is the reduced first Eilenberg subcomplex of the
singular complex $SK$ \cite[Definition~8.3 and Theorem~8.4]{May}.
In particular, as explained below,
\[
 |P_\bullet|\simeq\Omega K,\qquad
 \pi_1P_\bullet\cong\pi_2K.
\]
We retain this model of the presentation; we do not add further
nondegenerate generators to turn it into a resolution of $G$.
No nilpotency assumption on $G$ or $K$ is made here.

Recall that the ``Kan loop group functor'' $G$ assigns to each reduced simplicial set $X$ a free simplicial group $GX$ subject to the ``principal twisted cartesian product'' \cite[Theorem 26.6]{May}
$GX\rightarrow GX \times_{\tau} X \xrightarrow{p} X,$
where $|GX \times_{\tau} X|$ is a contractible space, $p$ is a Kan fibration and therefore $GX$ is a simplicial ``loop space'' for $X$.
By \cite[Theorem 10.10]{GJ}, the geometric realization $|p|$ of the Kan fibration $p$ is a Serre fibration. Then by \cite[Theorem 10.9]{GJ} $|GX|$ is the fiber of 
$|GX|\rightarrow |GX \times_{\tau} X| \xrightarrow{|p|} |X|$
and, therefore, by \cite[Proposition 4.66]{Hut} there is a weak homotopy equivalence $\beta:|GS_1K(X|R)|\simeq \Omega K(X|R)$.
Since $K(X|R)$ is a $CW$-complex, by \cite[Corollary 2]{Milnor} $\Omega K(X|R)$ is homotopy equivalent to a $CW$-complex.
And, as geometric realization is a CW-complex \cite[Theorem 14.1]{May}, by Whitehead's theorem, $\beta$ is a homotopy equivalence. Finally, we have a homotopy equivalence $|B_{K(X|R)}|\simeq \Omega K(X|R)$ and therefore $B_{K(X|R)}$ is a combinatorial model of based loops on $K(X|R)$.


%The free crossed module $(H/P,\Phi(X),\widetilde{\theta})$ of the presentation $(X|R)$ also arises as a second step on a way of constructing a free simplicial resolution from the given presentation. 


%One has $N(F^{(1)}_{\bullet})_1=Ker\; d_0\cong \Phi(Y\times \Phi(X))$  \cite[Proposition 1]{BH1} and we get the so called \textbf{free pre-crossed module of a
%discrete presentation} \eqref{2} on the set $Y$ as a pair
%$\Phi(Y\times \Phi(X))\xrightarrow{d_1} \Phi(X),$
%where the right action of $\Phi(X)$ is given by $a^u=s_0^0(u)^{-1}as_0^0(u), u \in \Phi(X), a \in \Phi(Y \times \Phi(X))$.

%From now on we omit the superscripts in $d^1_i$ and $s^0_0$ since they are already obvious in the 2-dimensional case.
%\begin{lemma} \label{l3.3}
%Let we are given a simplicial presentation \eqref{2}, then there is an isomorphism of crossed modules
%$\left(H/P, \Phi(X),
%\theta\right)\cong \left(\frac{Ker\; d_0}{[Ker\; d_0, Ker\; d_1]},\Phi(X),\overline{d_1}\right),$
%arising from the coincidence of the Peiffer subgroup $P$ and $[Ker\; d_0, Ker\; d_1]$ in $Ker\; d_0$.
%\end{lemma}

%By Brown-Loday Lemma \cite[5.7]{BL} and \cite{MP}, in any simplicial group degenerate in dimension two, in particular in \eqref{2}, the subgroups $[Ker\; d_0, Ker\; d_1]$ and $Im\;\overline{d_2}$ coincide,
%where $\overline{d_2}:N(F^{(1)}_{\bullet})_2\rightarrow N(F^{(1)}_{\bullet})_1$ - the restriction of $d^2_2$ onto $N(F^{(1)}_{\bullet})_2$.
%Let $C=Ker\; d_0/Im\;\overline{d_2}$, it is easy to check that $(C, \Phi(X),\overline{d_1})$ is a crossed module and we recollect all the above in the following
%\begin{proposition} \label{l3}
%Assume we are given a simplicial presentation \eqref{2}, then there is an isomorphism of crossed modules
%$\left(C, \Phi(X),\overline{d_1}\right)\cong\left(\frac{Ker\; d_0}{[Ker\; d_0, Ker\; d_1]},\Phi(X),\overline{d_1}\right),$
%arising from the coincidence of the subgroups $[Ker\; d_0, Ker\; d_1]=Im\;\overline{d_2}$ in $Ker\; d_0$ and hence the isomorphism
%$\pi_1F^{(1)}_{\bullet}\cong \frac{Ker\; d_0\cap Ker\;d_1}{[Ker\;d_0,Ker\; d_1]}\cong Ker\;\overline{d_1},$
%which allow us to identify $\pi_2K(X\mid R)$ with $\pi_1F^{(1)}_{\bullet}$.
%$ (C,F(X),\overline{d_1})=$
%\end{proposition}

A pro-$ p $-group is a group isomorphic to the inverse limit of finite $ p $-groups. This is a topological group (with inverse limit topology) which is compact and totally disconnected.
For such groups one has a presentation theory similar in many aspects to the combinatorial theory of discrete groups.

We next recall a classical arithmetic square under a nilpotency
hypothesis. This hypothesis applies only to the auxiliary statement
below; later it will be used at the nilpotent stages of the tower.
The general square in Proposition~\ref{general-free-arithmetic} does
not assume that the original space or simplicial group is nilpotent.
For a connected nilpotent space $X$ of finite type, the arithmetic square
of Bousfield and Kan \cite[Chapter~VI, Lemma~8.1]{BK} is a homotopy
pullback.  Write $\rho_0^T:T\to\mathbb Q_\infty T$ for the natural
rational-completion map and
$\kappa=\prod_{p\in\pi}\kappa_p$, where $\pi$ is the set of all primes.
Then the square is
\begin{equation}\label{a2}
\xymatrix{
 X\ar[rrr]^{\kappa}\ar[d]_{\rho_0^X} &&&
 \prod_{p\in\pi}(\mathbb F_p)_\infty X\ar[d]^{\rho_0^{\prod_p(\mathbb F_p)_\infty X}}\\
 \mathbb Q_\infty X\ar[rrr]^{\mathbb Q_\infty(\kappa)} &&&
 \mathbb Q_\infty\bigl(\prod_{p\in\pi}(\mathbb F_p)_\infty X\bigr).}
\end{equation}
Here the vertical arrows are maps, whereas $\mathbb Q_\infty$ denotes a
functor; the bottom arrow is its value on $\kappa$.  On nilpotent spaces,
rational completion agrees up to homotopy with rationalization
\cite[Chapter~V, Section~4.3]{BK}.  The corresponding group constructions
are rational prounipotent completion and pro-$p$ completion.  For the
interpretation of pro-$p$ groups as prounipotent group schemes over
$\mathbb F_p$, see \cite[Section~4.1]{Mikh2023}.
\begin{definition}\label{d4}\cite[2]{Mikh2023}
Let us fix a group $G$ and a field $k$ of zero characteristics, define the prounipotent completion of $G$ as the following universal diagram, in which $\rho$ is a
Zariski dense homomorphism from $G$ to the group of $k$-points
of a prounipotent affine group scheme $G_k^{\wedge}$:
$$\xymatrix @R=0.5cm{
                &         G^{\wedge}_k(k)  \ar[dd]^{\tau}     \\
 G \ar[ur]^{\rho} \ar[dr]_{\chi}                 \\
                &         H(k)              }$$
We require that for each Zariski dense homomorphism
$\chi$ there exists a unique homomorphism $\tau$ of prounipotent groups making the diagram commutative.
When $k=\mathbb{F}_p$, we define the prounipotent completion of $G$ as the pro-$p$-completion, which is a pro-$p$-group $G^{\wedge}_p$ obeys the same universal diagram, where $H(k)$ is a finite $p$-group ($G^{\wedge}_p\cong\varprojlim_{|G/U_{\lambda}|=p^{k_{\lambda}}} G/U_{\lambda}$).
\end{definition}

Bousfield and Kan show \cite[Proposition 4.1]{BK} that for each simplicial reduced space $X$ of finite type one can construct $R_{\infty}$-completion up to homotopy ($R=\mathbb{Q},\mathbb{F}_p$) in three steps:
\begin{description}
\item[i]{Replace $X$ with the so-called ``Kan's loop group'' $GX$ \cite[p.276]{GJ} - a simplicial group that has the homotopy type of ``loops on $X$'';}
\item[ii]{Apply the $R$-prounipotent completions functor dimension-wise, that is, we obtain a simplicial prounipotent group $(GX)^{\wedge}_R$;}
\item[iii]{Take the classifying space $\overline{W}(GX)^{\wedge}_R$ of $(GX)^{\wedge}_R$.}
\end{description}

The Kan loop-group functor $G$ and the classifying-space functor
$\overline W$ form an adjoint pair inducing the equivalence of the
corresponding homotopy categories \cite[Chapter~V, Proposition~6.3]{GJ}.
Let $P_\bullet$ be the free simplicial presentation associated with
$K(X\mid Y)$.  The relation between presentation homotopy groups and
simplicial-group homotopy groups is
\[
 \pi_n K(X\mid Y)\cong\pi_n(\overline WP_\bullet)
 \cong\pi_{n-1}(P_\bullet)\qquad(n\geq2).
\]
In particular, the second homotopy group of the presentation is
$\pi_2K(X\mid Y)\cong\pi_1(P_\bullet)$.  The index shift concerns
$\overline WP_\bullet$, not the realization $|P_\bullet|$, whose homotopy
groups agree with those of $P_\bullet$ without a shift.  The same
convention defines the second homotopy group of a prounipotent
presentation \cite[Section~4.5]{Mikh2023}.

\subsection{Outline of the argument}\label{intro:strategy}
For a finite presentation, write
$\widehat P_{R,\bullet}=(P_\bullet)^{\wedge}_R$ for the degreewise
rational or pro-$p$ completion of its simplicial model.  Then
\[
 R_\infty K(X\mid Y)\simeq\overline W((P_\bullet)^{\wedge}_R),
 \qquad R=\mathbb Q,\mathbb F_p.
\]
This model allows the three constructions described at the beginning
of the introduction to be carried out on the same simplicial object.
We outline their roles and the order in which they enter the argument.

For a finite subcontractible presentation, the ambient contractible
complex supplies a primitive independent family of abelianized
relators (Lemma~\ref{lem:relator-independence}).  After rational or
pro-$p$ completion, these relators extend to a free topological basis
$Y_1\sqcup Z$, and the quotient group is $F_R(Z)$, the free
prounipotent group read on rational points or the free pro-$p$ group,
respectively.  The canonical
relator-basis condition of Definition~\ref{d7} makes the cohomological
asphericity criterion applicable to this presentation, including the
non-minimal case.  The Brown--Loday identity identifies the resulting
low-dimensional obstruction with $\pi_1\widehat P_{R,\bullet}$.
To compute the higher homotopy groups we use the filtration, whose
initial input can be read directly from the chosen words.

Put $A_\bullet=\operatorname{Ab}_R(\widehat P_{R,\bullet})$.
Its normalized complex is
\[
 0\longrightarrow R^{Y_1}
 \xrightarrow{\,e_y\mapsto\operatorname{ab}_R(r_y)\,}
 R^{X_1}\longrightarrow0,
\]
in degrees one and zero.  The differential is a split injection, so
$A_\bullet$ is simplicially homotopy equivalent to the constant
module $R^Z$.  The closed lower central series in characteristic zero
and the Zassenhaus filtration in characteristic $p$ have graded
pieces $L_s^R(A_\bullet)$, where $L_s^R$ is the ordinary or restricted
free Lie power.  Thus the associated Adams--Quillen spectral sequences
have
\[
 E^1_{n,s}=\pi_nL_s^R(A_\bullet)=0
 \qquad(n>0,\ s\geq1).
\]
The mechanism is that a dimensionwise Lie functor preserves the
simplicial homotopy to the constant module.  Theorem~\ref{main}
then verifies convergence on the successive quotient fibrations:
the graded vanishing gives vanishing at every finite filtration
stage, and Milnor's exact sequence carries it to the complete
simplicial group.  This separates the computation of the spectral
sequence from the justification of its limiting conclusion.

Comparison of resolutions gives a further structural description of
this computation.  Once the positive homotopy groups have been shown
to vanish, the augmentation
$\widehat P_{R,\bullet}\to F_R(Z)$ is a free simplicial resolution.
Keune's comparison theorem lifts maps between the augmented groups
to maps of resolutions, uniquely up to simplicial homotopy.  In the
prounipotent category, Proposition~\ref{prop:keune-finite-type} and
Corollary~\ref{cor:prounipotent-model-comparison} therefore realize
\[
 \widehat P_{\mathbb Q,\bullet}\simeq
 \operatorname{Const}(F_{\mathbb Q}(Z))
\]
by morphisms and homotopies in that category.  The inductive step
lifts compatible full boundaries while preserving the values
already prescribed on degeneracies; homotopy uniqueness follows
from a relative extension across a simplicial cylinder.  The
continuous lifting argument, also for infinite convergent bases and
in the pro-$p$ category, is developed in
\cite[Proposition~2.3 and Theorem~4.2]{MikhContinuous}.
Because the comparison maps and homotopies are morphisms in the
completed category, degreewise abelianization and filtration
quotients can be applied to them.  This explains how the comparison
of resolutions relates the models used in the spectral calculation.
In this application the resolving property is established before
the comparison with the constant resolution is invoked.

The arithmetic square assembles the completions through their
common nilpotent quotients.  For any free discrete simplicial group
$P_\bullet$ of finite type, put
$N_{s,\bullet}=P_\bullet/\gamma_sP_\bullet$.  In each degree these
are finitely generated free nilpotent groups.  Their arithmetic
squares have an adelic entry given by a restricted product of
$p$-adic unipotent groups relative to their pro-$p$ subgroups
(Proposition~\ref{asnfn}).  Rational points map diagonally into this
entry.  Arithmetic factorization makes the corresponding square
of classifying spaces a connected homotopy pullback.  Taking the
homotopy inverse limit of these squares gives
\[
 \xymatrix{
 \widehat P_{\mathbb Z,\bullet}\ar[r]\ar[d] &
 \prod_p\widehat P_{p,\bullet}\ar[d]\\
 \widehat P_{\mathbb Q,\bullet}\ar[r] & \mathcal A(P)_\bullet,}
\]
where $\widehat P_{\mathbb Z,\bullet}=\varprojlim_sN_{s,\bullet}$
and the other completed groups are formed degreewise.  The
lower-right entry is the explicit simplicial group
\[
 \mathcal A(P)_\bullet=
 \varprojlim_s\varinjlim_{S\subset\pi,\ S\text{ finite}}
 \left(
 \prod_{p\in S}((N_{s,\bullet})_p^{\wedge})_0
 \times\prod_{p\notin S}(N_{s,\bullet})_p^{\wedge}
 \right).
\]
The finite exceptional set of primes may depend on the nilpotent
stage $s$.  This feature accounts for the order of the two limits
and is examined in Section~\ref{adelic-explicit}.  The square is a
homotopy pullback by Proposition~\ref{general-free-arithmetic};
its construction uses nilpotency at the quotient stages and imposes
no nilpotency condition on $\pi_0P_\bullet$.  Applying the derived
classifying-space functor identifies its upper-left entry with the
integral Bousfield--Kan completion of $\overline WP_\bullet$.

For a subcontractible presentation the same complementary words
can be chosen at every nilpotent stage.  They give compatible
contractions
\[
 P_\bullet/\gamma_sP_\bullet\simeq
 \operatorname{Const}\bigl(\Phi(Z)/\gamma_s\Phi(Z)\bigr).
\]
Consequently all four entries of the arithmetic square can be
identified by taking limits of explicit groups.  Theorem~\ref{c4}
gives, in particular,
\[
 R_\infty K(X_1\mid Y_1)\simeq K(F_R(Z),1),
 \qquad R=\mathbb Q,\mathbb F_p,\mathbb Z,
\]
with $F_{\mathbb Z}(Z)=\varprojlim_s\Phi(Z)/\gamma_s\Phi(Z)$.
Here the vanishing of the higher homotopy groups and the
surjectivity of the group towers make passage to the limit
explicit.  Subcontractible presentations thus exhibit both the
spectral and resolution-theoretic computation of the field
completions and their assembly in a noncommutative arithmetic square.

%Как показали Беррик-Хилман, для любого стягиваемого копредставления его конечное подкопредставление асферично тогда и только тогда, когда верна гипотеза асферичности Уайтхеда.
%При этом, как известно, если гипотеза Уайтхеда не верна, то накрытие, соответствующее радикалу Адамса, является ацикличным 2-комплексом.
%В начале 70-х Дрор показал, как можно исследовать ациклические пространства с помощью ациклических разложений и их алгебраических инвариантов.
%Для подстягиваемых копредставлений удобно использовать относительные ациклические разложения, которые строятся с использованием целочисленного пополнения Боуфилда-Кана. 
%Мы показываем, что целочисленное пополнение Боуфилда-Кана конечного подстягиваемого копредставления асферично и проводим вычисления в его разложении Дрора.
%Quillen’s plus construction applied to a space $X$, yields a
%universal map $\kappa : X \rightarrow X^+$, which is characterized by the fact that it quotients out
%the maximal perfect subgroup of $\pi_1X$ and induces isomorphisms in all homology theories
%(including homology with twisted coefficients). In general, a map between connected
%spaces satisfies this homological condition if and only if its homotopy fiber is acyclic. We denote the homotopy fiber of $\kappa : X \rightarrow X^+$ by $AX$.

%The acyclic space functor $AX$ was constructed by Dror using the Postnikov decomposition of $X$.

%It turns out that homotopy group of $\pi_nX$ and $\pi_nX^+$ are related 

%\begin{theorem}\label{nt1}

%Moreover, there is an epimorphism $ \pi_{n+2}\Phi\cong \pi_{n+1}F\rightarrow H_{n+1}F\cong H_2(\widetilde{G}; I[\widetilde{G}] \otimes_{\widetilde{G}} \pi_nX)$.
%Here $\widetilde{G}$ is the universal central extension of the maximal perfect subgroup $G$ of $\pi_1X$, and
%$I[\widetilde{G}]$ is the augmentation ideal of the integral group ring of $\widetilde{G}$.
%\end{theorem}

%It is useful to know $X^+$ of a subaspherical presentation. We demonstrate how it works with the following 

%\begin{theorem}\label{nt2}
%Let $X$ be a subcontractible presentation, then 
%$X^+$ is homotopy equivalent to the covering of the Bousfield-Kan conmpletion $\mathbb{Z}_{\infty}  X$ corresponding to the $\pi_1 X/\mathcal{P}\pi_1X$, where $\mathcal{P}\pi_1X$ is the maximal perfect subgroup of $\pi_1X$.
%$AX\simeq X_{P\pi_1X}$, where $X_{P\pi_1X}$ is a covering corresponding to $\mathcal{P}\pi_1X$. Let $p:X_{P\pi_1X}\rightarrow X$ be the covering, then there is a homotopy equivalence $X_{P\pi_1X}\simeq AX$.
%\end{theorem}
%\begin{proof} 
%this follows from the fact, that $\mathcal{P}\pi_1X$ coincide with the intersection of the lower central series of $\Gamma$ and from the universal property of the plus-construction.

%It is proven in an unpublished paper of J.-C. Haussmann that $X^+$ is an aspherical complex and therefore the long exact sequence of homotopy groups of the homotopy fibration $AX\rightarrow X \rightarrow X^+$ shows that 
%$\pi_iAX\cong \pi_iX (i\geq 2)$.
%By Emmanuel Dror that the homotopy fiber map $a:AX\rightarrow X$ is, up to homotopy, universal for maps of acyclic complexes into $X$ and therefore $p$ is fractured into a composition
%$a\cdot b$, where $b: X_{P\pi_1X}\rightarrow AX$ is the corresponding universal map. 
%The isomorphisms $(p_i)_*: \pi_iX_{P\pi_1X}\rightarrow \pi_iX (i\geq 2)$ and  $p_1: \pi_1X_{P\pi_1X}\rightarrow P\pi_1X\subset\pi_1X$ and Whitehead theorem imply the required homotopy equivalence.
%\end{proof}




%\includegraphics[scale=1]{bullejos faro1.jpg}

%\includegraphics[scale=1]{bullejos faro.jpg}



%If $K$ is not aspherical and $G=\pi_1(K)$, then the Adams radical $\mathcal{A}(G)=\cap\{ N\unlhd G : G/N \text{ is conservative}\}$ is not trivial and the corresponding regular covering $K_{\mathcal{A}(G)}\xrightarrow{\tau } K$ is acyclic (i.e. $H_1(K_{\mathcal{A}(G)})=H_2(K_{\mathcal{A}(G)})=0$) \cite[X, Theorem 6.2]{Bog}.
%In early 70's E.  Dror used acyclic Moore-Postnikov towers and showed how to analyze an acyclic space using the so called relative acyclic decomposition.
%$\mathbb{Z}_{\infty}  K$ is necessary ingredient when studing acyclic phenomena using the relative acyclic decomposition, since the latter is defined using pull-backs $\{K_k\}$ of $\mathbb{Z}_{\infty}  K$ and the Moore-Postnikov tower $\{P_k\}$ of $K$ \cite[4,5,8]{Dror73}

%$$\xymatrix{K_{\mathcal{A}(G)}\ar[r]^{\tau}& K \ar[rd]^{g_k}\ar[rdd]_{f_k}\ar[rrd]^{i}&&\\
%&&K_k \ar[r]^{\overline{i_k}}\ar[d]^{\overline{f_k}}  & \mathbb{Z}_{\infty}  K\ar[d]^{\mathbb{Z}_{\infty}f_k} \\
%&& P_k \ar[r]^{i_k} & \mathbb{Z}_{\infty} P_k K}.$$

%The plus construction $K(X_1|Y_1)^+$ of a subcontractible presentation $K(X_1|Y_1)$ (arising by gluing two and three dimensional discs) coincides up to homotopy with the Bousfield-Kan fiberwise $\mathbb{Z}$-completion of $K(X_1|Y_1)$ \cite[Definition 3.1.1.5]{Dun}. The universal property of the plus construction imply that there is a natural map $\tau: S_1K(X_1|Y_1)^+ \rightarrow \mathbb{Z}_{\infty}S_1K(X_1|Y_1)$.
%By Corollary \ref{c5}, $\mathbb{Z}_{\infty}S_1K(X_1|Y_1)$ is aspherical and it is also known \cite[Theorem 1]{Hau} (see also \cite[p. 122]{Gil}) that $S_1K(X_1|Y_1)^+$ is also aspherical.
%As sensitive information regarding the ``conservative radical" of $\pi_1K(X_1|Y_1)$ is contained in the 3-cells of the plus construction, we insist that study of $\tau$ definitely deserves further attention.

%\begin{theorem}{\cite[Theorem 4.9]{Mikh2023}}\label{000}
%Let $(X|Y_1)$ be a simplicial prounipotent  subpresentation of a simplicial aspherical prounipotent presentation $(X|Y)$  (if  $char(k)>0$ we assume that $(X|Y)$ is a simplicial pro-$p$-presentation), then $(X|Y_1)$ is also aspherical.
%\end{theorem}
%D. Anick's paper \cite{An} contain a simply connected version of the rational J.H.C Whitehead's problem i.e. it is assumed that the fundamental group is trivial: ``In rational homotopy we generally consider simply connected spaces only and tensor all homotopy groups with $\mathbb{Q}$''. At the time of writing, it was known, that the rational homotopy theory generalizes, at least, to the case of non-simply-connected nilpotent spaces \cite{BG} and therefore the restriction more reflects the lack of possibility of extending the methods of \cite{An} to a non-simply connected case.
Section~\ref{s1.0} develops the nilpotent arithmetic square and the
homotopy-limit constructions used to assemble its stages.
Section~\ref{s5} establishes independence of the relators, proves
the field-completion theorem by the filtration argument, and gives
the comparison of free simplicial prounipotent resolutions.
Section~\ref{s6} constructs the general noncommutative arithmetic
square, describes its adelic entry, and carries out the integral
completion calculation for subcontractible presentations.
%In Subsection \ref{s7} we propose an explanation of Serre's philosophy that homotopy types of pro-$p$-presentations determine the structure of discrete presentations in problems of (co)homological nature.

\section{Arithmetic squares for nilpotent spaces and their homotopy (co)limits}\label{s1.0}
\setcounter{theorem}{0}
We start with the "arithmetic square" of abelian groups \cite[Proposition 1.18]{Su1}
\begin{proposition}\label{su}
There is the following diagram of abelian groups
$$\xymatrix{\mathbb{Z}  \ar[rr]^{\prod_{p\in\pi} \kappa_p}\ar[d]_{\otimes \mathbb{Q}} && \prod_{p\in\pi}  \mathbb{Z}_p\ar[d]^{\otimes\mathbb{Q}} \\
\mathbb{Q}\ar[rr]^{\prod_{p\in\pi} \kappa_p} &&\widetilde{\prod}_{p\in\pi} \mathbb{Q}_p,}$$
where $\pi$ is the set of prime numbers, $\kappa_p$ is the $p$-adic
completion map, and $\widetilde{\prod}_{p\in\pi}\mathbb Q_p$ is the
restricted product: all but finitely many components belong to
$\mathbb Z_p$.
\end{proposition}
\begin{proof}
It is necessary to check the exactness of the exact sequence as follows
$$0\rightarrow \mathbb{Z}\rightarrow \prod_{p\in\pi}  \mathbb{Z}_p\oplus\mathbb{Q}  \xrightarrow{\otimes \mathbb{Q}-\prod_{p\in\pi} \kappa_p} \widetilde{\prod}_{p\in\pi} \mathbb{Q}_p\rightarrow 0.$$
First we note that for any $p$-adic number $r_p\in\mathbb{Q}_p$ there is a minimal $n\in\mathbb{N}$ such that $r_p=1/p^n\cdot \xi_p,$ where $\xi\in\mathbb{Z}_p$. Since $\mathbb{Z}_p$ is the $p$-adic completion of $\mathbb{Z}$ it follows that elements of $\mathbb{Z}_p$ can be performed as series of the form $\sum_{i=0}^{\infty} a_i\cdot p^i$, where $a_i$ are non-negative integers not exceeding $p-1$. It follows that for each $p$-adic number $r_p$ there is a rational number $m/p^n$, such that $r_p-m/p^n$ is a $p$-adic integer and therefore the sequence is right exact. It is also left exact since a rational number is a $p$-adic integer if and only if it is an integer.
\end{proof}
The $p$-adic completion and the rationalization can be extended to the category of groups.

Let $G$ be a group and let $p$ be a prime.  If $I$ is the augmentation
ideal of $\mathbb F_p[G]$, the Zassenhaus filtration is
\[
 \mathcal M_s(G)=\{g\in G\mid g-1\in I^s\}.
\]
Equivalently, the dimension-subgroup formula gives
\[
 \mathcal M_s(G)=\prod_{ip^j\geq s}\gamma_i(G)^{p^j},
 \qquad i\geq1,\quad j\geq0;
\]
see \cite{Rec,Cur} and \cite[Introduction and Section~8]{ChapmanEfrat2016}.
For pro-$p$ groups we use the completed group algebra, closed powers
of its augmentation ideal, and the closure of the displayed product.
Thus the augmentation-ideal definition here and the Zassenhaus
filtration used in the field-completion proof describe the same
filtration.  To avoid a terminological ambiguity, we distinguish it
from the lower $p$-central series
\[
 P_1G=G,\qquad P_{s+1}G=(P_sG)^p[P_sG,G],\qquad
 P_sG=\prod_{i=1}^s\gamma_i(G)^{p^{s-i}}.
\]
The latter product formula is the one discussed in
\cite[Introduction and Corollary~7.5]{ChapmanEfrat2016}, with reference
to \cite[Proposition~3.8.6]{NSW}.  These filtrations need not agree
term by term: for $G=\mathbb Z$ and $p=2$,
$\mathcal M_4(G)=4\mathbb Z$, whereas $P_4G=8\mathbb Z$.

Now we can define the $p$-adic completion of a group $G$ by the rule as follows $$G^{\wedge}_p = \varprojlim_s G/\mathcal{M}_s.$$
in the special case, when $G$ is finitely generated, then  $G^{\wedge}_p$ is a pro-$p$-completion of $G$ and thus, if $G$ is abelian, then $G^{\wedge}_p\cong G\otimes \mathbb{Z}_p$.

For a nilpotent group $G$, its rationalization $G_0$ is its Malcev
completion: the universal map to a uniquely divisible nilpotent group.
For an arbitrary group we instead distinguish the rational pronilpotent
completion
\[
 \widehat G_{\mathbb Q}^{\mathrm{nil}}
 =\varprojlim_s(G/\gamma_sG)_0.
\]
Let $J$ be the augmentation ideal of $\mathbb Q[G]$.  Quillen's
complete Hopf-algebra description identifies this completion with
the group-like elements of
\[
 \widehat{\mathbb Q[G]}=\varprojlim_m\mathbb Q[G]/J^m;
\]
see \cite[Appendix~A, Section~3]{Qui5}.  The Hopf structure is understood
on the completion, with completed tensor products.  The individual
algebra quotients by $J^m$ are not being asserted to be Hopf quotients.
The augmentation filtration and the nilpotent description give
compatible models of the same group completion.  Indeed, every map
to a rational nilpotent group of class at most $s$ kills
$\gamma_{s+1}G$ and factors uniquely through $(G/\gamma_{s+1}G)_0$.
Consequently these rationalized nilpotent quotients suffice to compute
the universal inverse limit, and all the identifications are natural
in $G$.  In the free finitely generated case used below, the Magnus
expansion identifies $\widehat{\mathbb Q[G]}$ with noncommutative power
series; its primitive elements form the completed free Lie algebra,
and truncation by bracket length gives the same nilpotent tower.
This group completion should not be confused with localization in
the category of all groups.

We recall the localization terminology only in the category of nilpotent
groups, where it is needed for the arithmetic square
\cite[Definition~1.1]{HMR}.
\begin{definition}\label{hmr}
Let $P$ be a set of primes.  A nilpotent group $G$ is \emph{$P$-local}
if $g\mapsto g^q$ is bijective for every prime $q\notin P$.
A \emph{$P$-localization} of $G$ is a homomorphism
$\kappa_P:G\to G_P$ to a $P$-local nilpotent group such that
precomposition induces a bijection
\[
 \operatorname{Hom}(G_P,H)\xrightarrow{\ \sim\ }
 \operatorname{Hom}(G,H)
\]
for every $P$-local nilpotent group $H$.
\end{definition}
Localization and completion are different operations.  Already for the
additive group of integers,
\[
 \mathbb Z_{\{p\}}=\mathbb Z_{(p)}
   =\{a/b\in\mathbb Q:p\nmid b\},
 \qquad \mathbb Z_p^{\wedge}=\mathbb Z_p.
\]
The former inverts primes other than $p$; the latter is the inverse
limit of the quotients $\mathbb Z/p^m\mathbb Z$.  The first group is
countable and the second uncountable.  Rationalization is the special
case $P=\emptyset$, denoted $G_0$, and coincides with Malcev completion
for nilpotent groups.  At the space level, Bousfield--Kan rational
completion agrees with rationalization for nilpotent spaces, but this
identification cannot be assumed for arbitrary spaces
\cite[Sections~7--8]{IvanovRat2021}.  We therefore retain separate
notation for localization, pro-$p$ completion, group pronilpotent
completion, and Bousfield--Kan completion.

It is proved in \cite[Lemma 8.2, Ch. VI]{BK}, \cite[Theorem I.3.7]{HMR} (here the proof is very similar to that of Proposition \ref{su})that any finitely generated nilpotent group is a pullback of the "arithmetic square" as follows
\begin{theorem}\cite[Lemma 8.2, Ch. VI]{BK} \label{asn}
If $G$ is a finitely generated nilpotent group and $P$ is a set of primes, then $G_P$ is a pullback of the arithmetic square below.  In this diagram $G_p$ denotes the pro-$p$ completion $G_p^{\wedge}$, whereas $G_P$ denotes localization at $P$.
$$\xymatrix{G_P  \ar[rrr]^{\prod_{p\in P} \kappa_p}\ar[d]_{\kappa_0} &&& \prod_{p\in P}  G_p\ar[d]^{\kappa_0} \\
 G_0 \ar[rrr]^{\kappa_0(\prod_{p\in P} \kappa_p)} &&&(\prod_{p\in P} G_p)^{\wedge}_0,}$$
Moreover, every element $u\in ( \prod_{p\in P} G_p)_0$ can be expressed as $u=vw$, where $v$ (resp. $w$) is in the image of 
$G_0$ (resp. $\prod_{p\in P} G_p$). If we assume the set $P=\pi$, then we obtain a pull-back presentation of a finitely generated nilpotent group in the "arithmetic square".
\end{theorem}


By a $\lambda$-nilpotent group we mean a group $G$ with
$\gamma_{\lambda+1}G=1$, where $\gamma_1G=G$ and
$\gamma_{j+1}G=[G,\gamma_jG]$.  We now make the lower-right corner of
the arithmetic square explicit.  Finite generation, rather than merely
finite rank, is assumed throughout this statement.
\begin{proposition}\label{asnfn}
Let $G$ be a finitely generated nilpotent group.  Put $N_p=G_p^{\wedge}$,
$N=\prod_pN_p$, and let $U$ be the unipotent algebraic group over
$\mathbb Q$ associated with its Malcev completion, so that
$G_0=U(\mathbb Q)$.  Write $U_p=U\otimes_{\mathbb Q}\mathbb Q_p$ and
$j_p:N_p\to U_p(\mathbb Q_p)$ for the natural map.  Then:
\begin{description}
\item[(1)] The rationalization of the underlying nilpotent group $N_p$
is naturally $(N_p)_0\cong U_p(\mathbb Q_p)$.
\item[(2)] For finite sets $S$ of primes, put
\[
 A_S=\prod_{p\in S}U_p(\mathbb Q_p)\times\prod_{p\notin S}N_p,
 \qquad A=\varinjlim_{S}A_S.
\]
For $S\subseteq T$, the transition map applies $j_p$ in the newly added
coordinates $p\in T\setminus S$ and is the identity elsewhere.  There
is a natural isomorphism $N_0\cong A$.  Equivalently, with
$K_p=j_p(N_p)$,
\[
 A\cong\prod_p'\bigl(U_p(\mathbb Q_p),K_p\bigr),
\]
the group of tuples whose coordinates belong to $K_p$ for all but
finitely many primes.
\item[(3)] Let $\kappa:G\to N$ be the product of the completion maps,
and let $\delta=\kappa_0:G_0\to N_0\cong A$ be its rationalization.
The arithmetic square is the pullback
\begin{equation}\label{fn}
\xymatrix{
 G\ar[rr]^{\kappa}\ar[d]_{\kappa_0} &&N\ar[d]^{\kappa_0}\\
 G_0\ar[rr]^{\delta} &&A.}
\end{equation}
In restricted-product coordinates,
$\delta(u)=(u_p)_p$, where $u_p$ is the image of $u\in U(\mathbb Q)$
under scalar extension to $\mathbb Q_p$.  This is a diagonal map over
\emph{all} primes.  By contrast, the map of an individual factor
$N_p\to A$ has only its $p$-coordinate possibly nontrivial and factors
through $j_p$.
\item[(4)] If $G$ is a free nilpotent group on finitely many generators,
all four arrows in \eqref{fn} are injective.  Thus, inside $A$,
\[
 G\cong\delta(G_0)\cap\kappa_0(N).
\]
\end{description}
Here $U_p(\mathbb Q_p)$ is also denoted
$G_{\mathbb Q_p}^{\wedge}(\mathbb Q_p)$; the notation distinguishes the
algebraic group from its group of rational points.
\end{proposition}
\begin{proof}
(1) The Malcev group of a finitely generated nilpotent group is
finite-dimensional.  The comparison of unipotent completions identifies
the $p$-adic unipotent completion with $U_p$; see
\cite[Appendix~A, Theorem~A.6 and Lemma~A.7]{HM2003}.
The map $j_p$ has torsion kernel, and its image $K_p$ is a compact open
subgroup of $U_p(\mathbb Q_p)$.  These facts can also be seen by taking
Malcev coordinates on $G/\operatorname{Tor}(G)$: the pro-$p$ completion
replaces integral coordinates by $p$-adic integral coordinates.  The
finite torsion subgroup contributes only a finite $p$-group to the
kernel; compare \cite[Appendix~A, Theorem~A.3 and Lemma~A.5]{HM2003}.
Every $v\in U_p(\mathbb Q_p)$ has a power in $K_p$.  Indeed,
$v^{p^m}=\exp(p^m\log v)$ tends to the identity, hence lies in the open
subgroup $K_p$ for large $m$.  Since $U_p(\mathbb Q_p)$ is uniquely
divisible and nilpotent, the rationalization criterion
\cite[Appendix~A, Corollary~3.8]{Qui5} proves (1).

(2) The nilpotency classes of all $A_S$ are bounded by that of $G$.
For a pro-$p$ group the $n$th-power map is bijective when $(n,p)=1$:
this holds on every finite $p$-group and passes to the inverse limit.
Given $a\in A_S$ and $n\geq1$, enlarge $S$ by the primes dividing $n$.
In the enlarged coordinates take roots in the uniquely divisible groups
$U_p(\mathbb Q_p)$; outside them take the unique roots in $N_p$.
This proves existence of $n$th roots in $A$.  If two elements have the
same $n$th power in the colimit, their equality of powers holds at some
finite stage; enlarging that stage by the primes dividing $n$ proves
uniqueness in exactly the same way.  Thus $A$ is uniquely divisible.

Let $j:N=A_{\emptyset}\to A$ be the canonical map.  For $a\in A_S$,
part (1) supplies a positive integer $n$ such that every coordinate in
$S$ of $a^n$ lifts to $N_p$.  The remaining coordinates already lie in
$N_p$, so $a^n\in j(N)$.  If $b\in\ker j$, its image is trivial at
some finite stage $A_S$.  Hence $b_p=1$ outside $S$, while each $b_p$
inside $S$ belongs to the torsion kernel of $j_p$.  A common multiple
of their orders annihilates $b$.  Conversely, all torsion elements map
to the torsion-free group $A$.  The same rationalization criterion now
gives $N_0\cong A$.

The coordinate maps identify $A$ with the displayed restricted product.
Surjectivity follows by lifting every coordinate in $K_p$ outside a
finite exceptional set.  For injectivity, note that
$\operatorname{Tor}(G)$ is finite, so the kernels of $j_p$ are trivial
outside a fixed finite set of primes.  An element with all coordinates
trivial therefore becomes trivial after adjoining that finite set to
$S$.  This also explains why the transition maps need not be embeddings
when $G$ has torsion.

(3) Every $u\in G_0$ has a positive power $u^n$ in the image of $G$.
For $p\nmid n$, the unique $n$th root of the image of $u^n$ in $N_p$
maps to $u_p$.  Therefore $u_p\in K_p$ except possibly at the finitely
many primes dividing $n$, and $(u_p)_p$ belongs to the restricted
product.  Its map from $G_0$ agrees with $j\kappa$ on $G$, hence equals
$\delta$ by the universal property of rationalization.  The pullback
assertion is Theorem~\ref{asn} with all primes.  Notice that an embedding
into a single coordinate of $A$ would not give this diagonal map.

(4) Write $G=\Phi(X)/\gamma_{\lambda+1}\Phi(X)$ with $X$ finite.
The Hall basic commutators of weights at most $\lambda$ give unique
integral coordinates on $G$.  The same coordinates are in $\mathbb Z_p$
for $N_p$, in $\mathbb Q$ for $G_0$, and in $\mathbb Q_p$ for
$U_p(\mathbb Q_p)$.  The inclusions of these coefficient rings show that
$G\to G_0$, $G\to N_p$, and $j_p$ are injective.  The coordinatewise
map $N\to A$ and the diagonal map $G_0\to A$ are consequently
injective, proving (4).

For clarity, the associated-graded argument behind these coordinates
is as follows.  Witt's theorem gives
\[
 \operatorname{gr}_{\gamma}\Phi(X)\cong L_{\mathbb Z}(X).
\]
In characteristic zero, the Lie algebra of the free prounipotent group
is the \emph{completed} free Lie algebra
\[
 \operatorname{Lie}F_k(X)=\widehat L_k(X)
 =\prod_{m\geq1}L_{k,m}(X).
\]
Its degree filtration has associated graded
$\bigoplus_{m\geq1}L_{k,m}(X)=L_k(X)$.  The exponential and logarithm
identify the group with this completed Lie algebra equipped with the
Baker--Campbell--Hausdorff multiplication, not with its additive group.
On each successive lower-central factor the multiplication becomes additive.
Since each $L_{\mathbb Z,m}(X)$ has a Hall basis, extension of scalars gives
\[
 \operatorname{gr}_{\gamma,m}F_k(X)(k)
 \cong L_{k,m}(X)
 \cong L_{\mathbb Z,m}(X)\otimes_{\mathbb Z}k.
\]
Here lower-central subgroups on the prounipotent side are closed;
see \cite[Chapter~IV, Sections~6--7]{Se} and
\cite[Appendix~A, Section~3]{Qui5}.
For the free pro-$p$ group, the closed ordinary lower-central quotients
are likewise
\[
 \operatorname{gr}_{\gamma,m}F_p(X)
 \cong L_{\mathbb Z,m}(X)\otimes_{\mathbb Z}\mathbb Z_p
\]
\cite[Propositions~3.2.2 and~3.2.5]{ZR}.  This is the ordinary closed
lower central series, not the Zassenhaus filtration.  Truncating these
identifications at weight $\lambda$ yields the free nilpotent case
used above.
\end{proof}


Bousfield and Kan show \cite[Proposition 4.1]{BK} that for each simplicial reduced space $X$ of finite type one can construct $R_{\infty}$-completion up to homotopy (we meet $R=\mathbb{Z}, \mathbb{Q},\mathbb{F}_p$) in three steps:
\begin{description}
\item[i]{Replace $X$ with the so-called ``Kan's loop group'' $GX$ \cite[p.276]{GJ} - a simplicial group that has the homotopy type of ``loops on $X$'';}
\item[ii]{Apply the $R$-prounipotent completions functor dimension-wise, that is, we obtain a simplicial prounipotent group $(GX)^{\wedge}_R$;}
\item[iii]{Take the classifying space $\overline{W}(GX)^{\wedge}_R$ of $(GX)^{\wedge}_R$.}
\end{description}

\begin{theorem}\cite[Chapter~V, Section~3]{BK}\label{q}
If $X$ is pointed, connected and nilpotent, then
$X\to\mathbb Q_\infty X$ is a rational homology equivalence and
\[
 \pi_1(\mathbb Q_\infty X)\cong(\pi_1X)_0,\qquad
 \pi_n(\mathbb Q_\infty X)\cong\pi_n(X)\otimes\mathbb Q
 \quad(n\geq2).
\]
The first formula uses nilpotent group rationalization; it is not a
formula involving the tensor product of a nonabelian group.
\end{theorem}
\begin{theorem}\cite[Chapter~VI, Section~5]{BK}\label{p}
Let $X$ be pointed, connected and nilpotent, with $\pi_1X$ finitely
generated and $\pi_nX$ finitely generated abelian for all $n\geq2$.
Then $X\to(\mathbb F_p)_\infty X$ is a mod-$p$ homology equivalence and
\[
 \pi_1((\mathbb F_p)_\infty X)\cong(\pi_1X)_p^{\wedge},\qquad
 \pi_n((\mathbb F_p)_\infty X)\cong\pi_n(X)\otimes\mathbb Z_p
 \quad(n\geq2).
\]
\end{theorem}
\noindent\emph{Why finite generation occurs in the second theorem.}
For a general abelian group $A$, derived $p$-completion involves
\[
 L_0^pA=\operatorname{Ext}^1_{\mathbb Z}(\mathbb Z/p^{\infty},A),
 \qquad L_1^pA=\operatorname{Hom}_{\mathbb Z}(\mathbb Z/p^{\infty},A).
\]
For example, for simply connected $X$ the general formula is
\[
 0\longrightarrow L_0^p(\pi_nX)
 \longrightarrow\pi_n((\mathbb F_p)_\infty X)
 \longrightarrow L_1^p(\pi_{n-1}X)\longrightarrow0
 \qquad(n\geq2);
\]
see \cite[Chapter~VI, Proposition~5.1]{BK}.  For finitely generated $A$,
$L_1^pA=0$ and $L_0^pA=A\otimes\mathbb Z_p$, which explains the simple
formula above.  Without this hypothesis additional homotopy can appear:
for $A=\mathbb Z/p^{\infty}$ one has $L_0^pA=0$ and
$L_1^pA\cong\mathbb Z_p$, so the completion of $K(A,2)$ is
$K(\mathbb Z_p,3)$.  Finite generation is thus a sufficient condition
for the tensor-product formula, not a requirement for defining
completion itself.  No corresponding finiteness hypothesis is needed
for the rational formula in Theorem~\ref{q}.

Countability alone does not suffice.  For the countable free abelian group
$A=\bigoplus_{j\geq1}\mathbb Z$, ordinary $p$-completion gives
\[
 \widehat A_p=\varprojlim_m A/p^mA
 =\left\{(a_j)\in\prod_{j\geq1}\mathbb Z_p:
       a_j\longrightarrow0\text{ $p$-adically}\right\}
 \supsetneq\bigoplus_{j\geq1}\mathbb Z_p=A\otimes\mathbb Z_p;
\]
the sequence $(p^j)_{j\geq1}$ witnesses strictness.
Conversely, under the hypotheses of Theorem~\ref{p}, writing
$\pi_nX\cong\mathbb Z^{r_n}\oplus T_n$ with $T_n$ finite yields
$\pi_n((\mathbb F_p)_\infty X)\cong
\mathbb Z_p^{r_n}\oplus(T_n)_{(p)}$ for $n\geq2$.
Thus these completed groups are finite when $r_n=0$ and have cardinality
$2^{\aleph_0}$ when $r_n>0$.  Finite type controls the completion
formula; it does not imply countability of the completed homotopy groups.

\subsection{Homotopy limits and colimits}\label{Holim} 
We assume that the reader is familiar with the notions of homotopy limits and colimits appropriate references to our taste are \cite{MV}, \cite{Hir2}.
However we remind some notations for instance.

We call $X$ a $I$-diagram of simplicial sets if $X$ is a functor $X: I\rightarrow SSets$ from the small category $I$ to the category $SSets$ of punctured simplicial sets
and denote the category of such diagrams by $Ssets^I$. We also need the notions of overcategory $(I\downarrow i)$ and of undercategory $(i\uparrow I)$, where $i\in Ob(I)$ \cite[2.1]{Hir2} and their classifying spaces (i.e. nerves of this categories $B(I\downarrow i)$, $B(i\uparrow I)$).
One have $X_0, X_1 : I^{op}\times I\rightarrow Ssets$ - two bifunctors $X_0(i,j)=hom(B(I\downarrow i), X(j))$ and $X_1(i,j):B(I\uparrow I)\times X(j)$, where $hom(A,B)$ is the simplicial mapping space - the internal $Hom$ functor in the category of simplicial sets (i.e. $hom_{SSets}(A,B)$ is also a simplicial set) given for simplicial sets $A,B$ by the formula $hom(A,B)_n=Hom_{SSets}(A\times \Delta^n,B)$.
Now the homotopy limit and colimit are the end and coend of $X_0$ and $X_1$ respectively
$$holim_I X=\int_I X_0\qquad\qquad \text{and}\qquad\qquad hocolim_I X=\int^I X_1.$$
Unravelling the end formula \cite[7.4.6, 7.4.11]{MV}, we view
\[
 \operatorname{holim}_I X\subseteq
 \prod_i\operatorname{hom}(B(I\downarrow i),X(i)).
\]
An $n$-simplex is a family $\xi=\{\xi_i\}_{i\in I}$ such that, for
every morphism $i\to j$, the following diagram commutes; the horizontal
maps are induced by this morphism and the identity of $\Delta^n$:
$$\xymatrix{B(I\downarrow i)\times \Delta^n  \ar[r]\ar[d]_{\xi_i} & B(I\downarrow j)\times \Delta^n\ar[d]_{\xi_j} \\
X(i)\ar[r] &X(j)}$$
commutes and therefore $holim_I X$ can be equivalently defined as a simplicial set of natural transformations of functors 
$holim_I X = hom_{SSets^I} (B(I\downarrow \:),X)$, where 
$hom_{SSets^I} (B(I\downarrow\:),X)_n=hom_{SSets^I}(B(I\downarrow\:)\times \Delta^n,X)$ is a simplicial function complex in the simplicial model category $SSets^I$ \cite[I.5]{GJ}.
Simplicial mapping space functor behaves cute with respect to geometric realization and total singular space functors \cite[I.1]{GJ}. In particular there is a natural isomorphism of the spaces of $n$-simplexes $$(Sing(X^K))_n\cong ((Sing\; X)^K)_n,$$
where $X^K$ is $CGHS$- the compactly generated Hausdorf space  of continuous functions from $|K|$- the geometric realization of $K$ to $X$ and $Sing\; X$ is the total singular simplicial space of $X$. Thus  in $SSets$ we have a natural isomorphism $Sing(X^K)\cong (Sing\;X)^K$.
As a consiquence the similar notions of homotopy limits and colimits for diagrams of topological spaced $T: I\rightarrow CGHS$ 
$$holim_I T=\int_I Hom_{CGHS}(|B(I\downarrow i)|, T(j)) \qquad\qquad hocolim_I T=\int^I |B(i\uparrow I)^{op}|\times T(j)$$ 
are compatible with simplicial ones. For instance, there are natural isomorphisms of topological spaces and simplicial sets respectively \cite[14.2]{Hir2} $|hocolim_I\;X|\cong hocolim|X|$, $Sing(holim\; X)\cong holim(Sing(X)))$ as well as natural weak equivalences for objectwise fibrant and cofibrant diagrams respectively $|holim\; X|\simeq holim|X|$ and $hocolim(Sing\; X)\simeq Sing(hocolim\; X)$.

We refer the reader to \cite{MV} and \cite{Hir2} regarding homotopy inverse (co)limits. Our basic examples:
\begin{example}[Homotopy pullback]
Let $I_{pb}$ be the category $\{1\}\to\{1,2\}\leftarrow\{2\}$,
and consider a diagram
\[
 X_1\xrightarrow{f_1}X_{12}\xleftarrow{f_2}X_2
\]
of Kan complexes.  For an arbitrary diagram of simplicial sets, first
apply a functorial objectwise fibrant replacement.

In the end formula for the homotopy limit,
$N(I_{pb}\downarrow\{1\})=N(I_{pb}\downarrow\{2\})=\Delta[0]$,
whereas $K=N(I_{pb}\downarrow\{1,2\})$ consists of two edges with a
common terminal vertex.  Although $|K|$ is homeomorphic to an interval,
$K$ is not isomorphic to $\Delta[1]$.  In particular, this homeomorphism
does not identify the simplicial mapping spaces $X_{12}^{K}$ and
$X_{12}^{\Delta[1]}$.  Likewise, the unit
$X\to\operatorname{Sing}|X|$ is a weak equivalence, not in general an
isomorphism of simplicial sets.

Instead, use the standard path-object factorization of the diagonal
\[
 X_{12}\xrightarrow{\simeq}X_{12}^{\Delta[1]}
 \xrightarrow{(\operatorname{ev}_0,\operatorname{ev}_1)}
 X_{12}\times X_{12}.
\]
The first map assigns the constant path, and the second is a Kan
fibration, since $X_{12}$ is fibrant and
$\partial\Delta[1]\hookrightarrow\Delta[1]$ is a cofibration; see
\cite[Chapter~I, Sections~5 and~9]{GJ}.  Here the cotensor uses the
unpointed interval; for a pointed target its base point is the constant
map to the base point.
Define
\[
 P=(X_1\times X_2)\times_{X_{12}\times X_{12}}
 X_{12}^{\Delta[1]}.
\]
Thus, degreewise, $P$ consists of triples $(x_1,x_2,\gamma)$ satisfying
$\operatorname{ev}_0\gamma=f_1(x_1)$ and
$\operatorname{ev}_1\gamma=f_2(x_2)$.  The path-object construction of
homotopy pullbacks gives a natural weak equivalence (or a natural zigzag
of weak equivalences, depending on the chosen homotopy-limit model)
\[
 P\simeq\operatorname{holim}_{I_{pb}}X;
\]
see \cite[Example~8.2.8]{MV}.  We henceforth use $P$ as our model for
$X_1\times^h_{X_{12}}X_2$.  In particular, the following square is an
ordinary pullback defining this model:
\begin{equation}\label{hp}
\xymatrix{
 P\ar[r]\ar[d]_{\operatorname{pr}} &
 X_{12}^{\Delta[1]}\ar[d]^{(\operatorname{ev}_0,\operatorname{ev}_1)}\\
 X_1\times X_2\ar[r]^{f_1\times f_2} & X_{12}\times X_{12}.}
\end{equation}
The projection $\operatorname{pr}$ is a Kan fibration.  If the spaces
and maps are pointed, its fiber over the pair of base points is the
based loop-space model $\Omega X_{12}$.

\begin{corollary}[Dyer-Roitberg]\label{dr}
For a diagram of pointed Kan complexes and pointed maps as above, there
is a fiber sequence
\[
 \Omega X_{12}\longrightarrow P
 \xrightarrow{\operatorname{pr}}X_1\times X_2.
\]
With compatible choices of base points, its homotopy exact sequence has,
for $n\geq2$, the segments
\[
 \cdots\longrightarrow\pi_{n+1}(X_{12})
 \longrightarrow\pi_n(P)
 \longrightarrow\pi_n(X_1)\times\pi_n(X_2)
 \xrightarrow{\ f_{1*}-f_{2*}\ }\pi_n(X_{12})
 \longrightarrow\pi_{n-1}(P)\longrightarrow\cdots,
\]
where the loop-space identification is chosen to give the displayed
sign convention.  The low-dimensional end is
\[
 \pi_1(P)\longrightarrow\pi_1(X_1)\times\pi_1(X_2)
 \longrightarrow\pi_1(X_{12})\longrightarrow\pi_0(P)
 \longrightarrow\pi_0(X_1)\times\pi_0(X_2).
\]
This end is interpreted in the usual sense of the homotopy exact
sequence of a fibration, with pointed sets and the fundamental-group
action where appropriate.  In particular, the map from
$\pi_1(X_1)\times\pi_1(X_2)$ to $\pi_1(X_{12})$ is not in general a
group homomorphism, and additive difference notation is not used there.
\end{corollary}
\begin{lemma}[Degreewise reconstruction of a homotopy pullback]
\label{degreewise-pullback}
In Corollary~\ref{dr}, suppose that $X_1,X_2,X_{12}$ are connected,
and put
\[
 \delta_n=f_{1*}-f_{2*}:\pi_nX_1\oplus\pi_nX_2\longrightarrow\pi_nX_{12}
 \qquad(n\geq2).
\]
There are natural short exact sequences
\[
 0\longrightarrow\operatorname{coker}\delta_{n+1}
 \longrightarrow\pi_nP\longrightarrow\ker\delta_n
 \longrightarrow0\qquad(n\geq2).
\]
Consequently, if $\delta_{n+1}$ is surjective, then
\[
 \pi_nP\cong\pi_nX_1\times_{\pi_nX_{12}}\pi_nX_2.
\]
If $\delta_2$ is surjective, the analogous assertion in degree one is
an isomorphism of groups
\[
 \pi_1P\cong\pi_1X_1\times_{\pi_1X_{12}}\pi_1X_2
\]
for the chosen component.  Moreover, $P$ is connected if the double
coset set
\[
 f_{1*}(\pi_1X_1)\backslash\pi_1X_{12}/f_{2*}(\pi_1X_2)
\]
is a singleton.  These identifications are natural in pointed maps
of the three-corner diagrams.
\end{lemma}
\begin{proof}
The short exact sequences are obtained by taking the indicated
segments of the homotopy exact sequence in Corollary~\ref{dr}.
Surjectivity kills the preceding boundary map, and the projection
from $P$ then identifies its homotopy group with the displayed kernel
or group pullback.  The components of $P$ are the orbits of
$\pi_1X_1\times\pi_1X_2$ on $\pi_1X_{12}$ under left and right
multiplication, giving the double coset description.  Naturality
follows from the path-object model and its homotopy exact sequence.
\end{proof}

Corollary~\ref{dr} and Lemma~\ref{degreewise-pullback}, together with
the algebraic arithmetic square already proved in
Theorem~\ref{asn}, give the arithmetic square for nilpotent spaces of
finite type.
\begin{theorem}\label{space-arithmetic}
Let $X$ be a connected nilpotent Kan complex of finite type.  Set
$X_0=\mathbb Q_\infty X$, $X_p=(\mathbb F_p)_\infty X$, and
$B=\prod_pX_p$.  Then the natural map
\[
 X\longrightarrow\overline X
 :=\operatorname{holim}(X_0\longrightarrow\mathbb Q_\infty B
 \longleftarrow B)
\]
is a weak equivalence, hence a homotopy equivalence after realization.
\end{theorem}
\begin{proof}
For nilpotent finite-type $X$, its fundamental group is finitely
generated and each higher homotopy group is finitely generated abelian.
The nilpotent actions on higher homotopy groups have finite filtrations;
completion preserves these filtrations, with a bound independent of $p$
in each degree.  Thus $B$ is nilpotent as well.  Put $G_n=\pi_nX$.
Theorems~\ref{q} and~\ref{p} give
\[
 \pi_nX_0=(G_n)_0,\quad
 \pi_nB=\prod_p(G_n)_p^{\wedge},\quad
 \pi_n(\mathbb Q_\infty B)=\left(\prod_p(G_n)_p^{\wedge}\right)_0,
\]
with ordinary abelian rationalization when $n\geq2$.
For every $n\geq2$ the arithmetic square gives the short exact sequence
\[
 0\longrightarrow G_n\longrightarrow
 (G_n)_0\times\prod_p(G_n)_p^{\wedge}
 \xrightarrow{\delta_n}\left(\prod_p(G_n)_p^{\wedge}\right)_0
 \longrightarrow0.
\]
Here exactness, including surjectivity of $\delta_n$, is precisely
Theorem~\ref{asn} applied to the finitely generated abelian group
$G_n$; it is not an additional assumption on the spaces.
Lemma~\ref{degreewise-pullback} therefore identifies
$\pi_n\overline X$ with $\ker\delta_n=G_n$.
In particular, the already proved factorization in degree $n+1$
eliminates the possible contribution of that degree.  The remaining
pullback is determined entirely by the arithmetic square of $G_n$.
For $n=1$, surjectivity of $\delta_2$ again gives injectivity into
$\pi_1X_0\times\pi_1B$; the image is the group-theoretic pullback,
which is $G_1$ by Theorem~\ref{asn}.
Finally, the factorization assertion of that theorem says that the two
images of $\pi_1X_0$ and $\pi_1B$ have a single double coset in
$\pi_1(\mathbb Q_\infty B)$.  Hence $\overline X$ is connected.
The natural map from $X$ induces precisely these identifications, and
is therefore a weak equivalence in all degrees.
All these isomorphisms are induced by the canonical comparison map;
no choices of arithmetic factorizations are used to define them.
They are therefore compatible with maps of spaces and, in particular,
with the transition maps of an inverse tower.
\end{proof}
\end{example}
\begin{example}[Homotopy inverse and direct limits]

The inverse category, $$I_s=\{... \rightarrow \bullet\rightarrow\bullet\rightarrow\bullet\}.$$ In this case $holim_{I_s} X$ is the so called homotopy inverse limit.
Dually, the direct category  $$D_s=\{... \leftarrow \bullet\leftarrow\bullet\leftarrow\bullet\}$$ here we call $hocolim_{D_s} X$ - the homotopy direct limit.

The categories $I_s$ and $D_s$ are particular examples of the so called categories of cofibrant and fibrant constants \cite[15.10]{Hir}. 
The corresponding categories  $SSets^{D_s}$ and  $SSets^{I_s}$ of $SSets$-diagrams have the canonical Reedy model categories structures, in which cofibrant diagrams $SSets^{D_s}$ are the direct diagrams, where each map is an embedding of simplicial sets and, dually, a fibrant diagram in $SSets^{I_s}$ is the inverse diagrams, whose initial object is Kan and whose transition maps are Kan fibrations. It follows from \cite[Theorem 19.9.1]{Hir} that for such (co)fibrant diagrams in $SSets^{I_s}$ ($SSets^{D_s}$) the natural maps $$lim \; X\rightarrow holim\; X\quad (hocolim\; X\rightarrow colim\; X)$$
are weak equivalences (and give homotopy equivalences after realization).
\end{example}
\begin{proposition}[Commuting homotopy limits and colimits]\label{hofubini}
For small categories $I,J$ and a diagram $X:I\times J\to SSets$, there
are natural weak equivalences
\[
 \operatorname{holim}_{I\times J}X\simeq
 \operatorname{holim}_I\operatorname{holim}_JX\simeq
 \operatorname{holim}_J\operatorname{holim}_IX,
\]
\[
 \operatorname{hocolim}_{I\times J}X\simeq
 \operatorname{hocolim}_I\operatorname{hocolim}_JX\simeq
 \operatorname{hocolim}_J\operatorname{hocolim}_IX.
\]
If $I$ is finite and acyclic and $J$ is filtered, the natural comparison
\[
 \operatorname{hocolim}_J\operatorname{holim}_IX
 \longrightarrow\operatorname{holim}_I\operatorname{hocolim}_JX
\]
is a weak equivalence, with fibrant replacements understood.  The last
assertion does not permit interchanging an infinite inverse tower with
a filtered homotopy colimit.
\end{proposition}
\begin{proof}
The first two assertions are the Fubini properties of homotopy
(co)limits; see \cite[Section~8.5.5]{MV}.
For the mixed assertion, a finite acyclic
category has only finitely many nondegenerate composable strings of
arrows.  Its homotopy limit is therefore a finite homotopy-limit
construction.  Filtered colimits of simplicial sets preserve weak
equivalences and commute with finite limits, and filtered colimits of
Kan complexes are Kan.  Applying these facts to a finite homotopy-limit
model proves the comparison; see also \cite[Theorem~14.17]{Hir2}.
\end{proof}

For completeness, homotopy invariance can be checked independently.
For objectwise fibrant diagrams, the Bousfield--Kan end model for
$\operatorname{holim}_I X$ is the mapping space of diagrams from the
cofibrant resolution $N(I\downarrow -)$ of the constant point to $X$.
The simplicial model-category axiom implies that this mapping-space
functor sends weak equivalences between fibrant diagrams to weak
equivalences.  For homotopy colimits, use the simplicial replacement,
whose terms are coproducts of the entries $X(i)$ indexed by composable
strings of arrows.  Objectwise weak equivalences give degreewise weak
equivalences of these bisimplicial sets and hence weak equivalences of
their diagonals.  In the pointed setting the corresponding constructions
use a disjoint base point on the indexing nerves and pointed coproducts.
These arguments justify passage to homotopy categories without assuming
that an arbitrary functor preserves weak equivalences.

The category $SSets^I$ is also a closed simplicial model category (the so called Bousfield-Kan model structure \cite[\S 8, Ch. XI]{BK}) with weak equivalences the maps $\phi:\mathcal{X}\rightarrow \mathcal{Y}\in SSets^I$such that $\phi_i:\mathcal{X}\rightarrow \mathcal{Y}\in SSets$ is a weak equivalence for every $i\in I$. The fibrations in $SSets^I$ are maps $\phi:\mathcal{X}\rightarrow \mathcal{Y}\in SSets^I$ 
such that $\phi_i:\mathcal X(i)\rightarrow\mathcal Y(i)$ is a fibration
for every $i\in I$; cofibrations are characterized by the left lifting
property with respect to objectwise trivial fibrations.  The simplicial
structure is obtained by viewing objects of $SSets^I$ as simplicial objects
in $Sets^I$ \cite[Section~5]{Hir2}.

\begin{definition} Assume that $\mathcal{C}$ is a model category and $F:\mathcal{C}\rightarrow \mathcal{D}$ is a functor. Consider pairs $(G,s)$ consisting of a functor $G: Ho(\mathcal{C})\rightarrow \mathcal{D}$ and natural transformation $s:G\gamma\rightarrow F$. A left derived functor for $F$ is a pair $(LF,t)$ of tis type which is universal from the left, in the sence that if $(G,s)$ is any such pair, then there exists a unique natural transformation $s':G\rightarrow LF$ such that the compositenatural transformation $$G\gamma\xrightarrow{s'\cdot \gamma} (LF)\gamma\xrightarrow{t} F$$ 
is the natural transformation $s$.
\end{definition}

The following proposition of Quillen \cite[]{Qui1} gives a sufficient condition of existence of left derived functors
\begin{proposition}
Let $\mathcal{C}$ is a model category  and $F:\mathcal{C}\rightarrow \mathcal{D}$ is a functor with the property that $F(f)$ is an isomorphism whenever $f$  is a weak equivalence between cofibrant objects 
in $\mathcal{C}$. Then the left derived functor $(LF,t)$ of $F$ exists, and for each cofibrant object $X$ of $\mathcal{C}$ the map
$$t_X:LF(X)\rightarrow F(X)$$
is an isomorphism.
\end{proposition}
We will also need Quillen's definition of a total derived functor.
\begin{definition}
Assume we are given $F:\mathcal{C}\rightarrow \mathcal{D}$ a functor between model categories. A total derived functor $LF$ for $F$ is a functor 
$$LF: Ho(\mathcal{C})\rightarrow Ho(\mathcal{D})$$
which is a left derived functor for the composite $\gamma_{\mathcal{D}}\cdot F: \mathcal{C}\rightarrow Ho(\mathcal{D})$.
Similarly, a total right derived functor $RF$ for $F$ is a functor $RF: Ho(\mathcal{C})\rightarrow Ho(\mathcal{D})$ which is a right derived functor for the composite $\gamma_{\mathcal{D}}\cdot F$.
\end{definition}

The constant-diagram (or diagonal) functor
\[
\Delta_I:SSets\longrightarrow SSets^I
\]
sends a simplicial set $K$ to the constant $I$-diagram with value $K$.
It fits into the familiar adjoint triple
\[
 \operatorname{colim}_I\dashv \Delta_I\dashv \operatorname{lim}_I.
\]

\begin{theorem}[Homotopy (co)limits as derived functors]
\label{thm:derived-holim-hocolim}
Let $I$ be a small category.  The Bousfield--Kan homotopy colimit and
homotopy limit induce functors
\[
 \operatorname{hocolim}_I:Ho(SSets^I)\longrightarrow Ho(SSets),
 \qquad
 \operatorname{holim}_I:Ho(SSets^I)\longrightarrow Ho(SSets).
\]
With respect to the adjoint triple
$\operatorname{colim}_I\dashv\Delta_I\dashv\operatorname{lim}_I$, there
are natural isomorphisms
\[
 \operatorname{hocolim}_I\cong \mathbf L\operatorname{colim}_I,
 \qquad
 \operatorname{holim}_I\cong \mathbf R\operatorname{lim}_I.
\]
Thus the homotopy colimit is the total left derived functor of the
ordinary colimit, whereas the homotopy limit is the total right derived
functor of the ordinary limit.
\end{theorem}

\begin{proof}
Equip $SSets^I$ with a model structure whose weak equivalences are
objectwise weak equivalences.  If $QX\xrightarrow{\sim}X$ is a cofibrant
replacement of an $I$-diagram, then the Bousfield--Kan coend formula is
naturally weakly equivalent to $\operatorname{colim}_I QX$.  Consequently
it sends weak equivalences between cofibrant diagrams to weak
equivalences and satisfies the universal property of the total left
derived functor of $\operatorname{colim}_I$.

Dually, if $X\xrightarrow{\sim}RX$ is a fibrant replacement in an
appropriate diagram model structure, then the Bousfield--Kan end formula
is naturally weakly equivalent to $\operatorname{lim}_I RX$.  It therefore
satisfies the universal property of the total right derived functor of
$\operatorname{lim}_I$.  These identifications are natural in $X$; see
\cite[Theorems~19.9.1 and~19.9.2]{Hir}.
\end{proof}

Notice that the theorem does not assert that homotopy (co)limits are
derived functors of $\Delta_I$.  Rather, $\Delta_I$ is the middle functor
in the two adjunctions, while the derived functors are taken from its left
and right adjoints.  Equivalently, on homotopy categories one obtains the
derived adjunctions
\[
 \operatorname{hocolim}_I\dashv Ho(\Delta_I)
 \dashv \operatorname{holim}_I.
\]

%From now on we omit the superscripts in $d^1_i$ and $s^0_0$ since they are already obvious in the 2-dimensional case.
%Recall that $\pi_1(F_{\bullet})$ coincides with $\pi_2(X|R)$ of the standard 2-complex of \eqref{eq1}.
%\begin{proof}
%We prove coincidence of $P$ and $[Ker\; d_0, Ker\; d_1]$ of two steps:

%1) using elementary substitutions we see that Peiffer commutators are representable in the form of
%commutators $[x^{-1}s_0d_1(x),y],$ where $x,y\in Ker\; d_0$;

%2) elements of $Ker\; d_1$ are exactly the elements  $x^{-1}s_0d_1(x),$ where $x\in Ker\; d_0.$

%The proof of CM 1) and CM 2) is straightforward.
%\end{proof}

%\subsection{The second homotopy module of a subcontractible presentation}\label{r1.1}

%\begin{theorem}\label{nth100}
%Assume $\overline{W}X$ be a finite type subcontractible presentation assotiated with a free simplicial group presentation $X$ of $\pi_0X$ then $\pi_2 \overline{W}X=0$.
%\end{theorem}
%\begin{proof}
%We consider the family of simplicial fibrations $\mathbb{Z}_s$ and their fibers $\chi_s, s\geq 2$ in the commutative diagram \cite[???]{BK} as follows 
%\begin{equation}\label{ne1}
%\xymatrix @R=0.5cm{
%  \chi_s \ar[dr]& & \\
%   \chi_{\infty}^c \ar[u]^{\mu_s} \ar[r]          &    \overline{W}X \ar[r]^{\mathbb{Z}_{\infty}} \ar[dr]_{\mathbb{Z}_s}   & \mathbb{Z}_{\infty}\overline{W}X\ar[d]^{\tau_s} \\
 %&              &      \overline{W}X/\Gamma_s X }
%\end{equation}
%arises in the homotopy equivalences $\mathbb{Z}_{\infty}\overline{W}X\simeq\varprojlim_s \mathbb{Z}_s \overline{W}X\simeq \varprojlim_s\overline{W}X/\Gamma_s X $ \cite[???]{BK}.

%By Theorem \ref{nth1} it follows that   $\overline{W}X/\Gamma_s X=0, s\geq 2$. 
%$X$ is a free simplicial group which is of finite rank in each dimension and degerated in dimensions greater than one. 
%Cosider the Milnor's exact sequence in homotopy groups assotiated with the inverse system 
%$0\rightarrow \varprojlim^1 \pi_3(\chi_s)\rightarrow \pi_2(\varprojlim_s\chi_s) \rightarrow\varprojlim_s \pi_2(\chi_s) \rightarrow 0$ of fibers $\chi_s$.

%Since $X_n, n\geq 0$ are finitely generated free discrete groups it follows $\varprojlim_s (\chi_s)_k\cong \cap_s (\chi_s)_k = 0, k\geq 0$ and therefore $\varprojlim^1 \pi_2\chi_s\cong \pi_1\varprojlim_s \chi_s \cong \Gamma$.

%But each $\chi_s$ is a 2-complex, being homotopy equivalent to a subcomplex of 2-complex and hence by \cite[]{Man09} $\pi_3\chi_s\cong S^2(\pi_2\chi_s)$, where $S^2((\pi_2\chi_s))= \{x\in \pi_2\chi_s\otimes \pi_2\chi_s| \sigma x=x \}$ and 
%$\sigma$ be the linear automorphism of $\pi_2\chi_s\otimes \pi_2\chi_s$ defined by $\sigma (a\otimes b)=b\otimes a$. 

%Now it is straightforward to see that $\varprojlim_s \pi_2(\chi_s) =0$, bur by Theorem \ref{nth1} we have an isomorphism $\pi_2 \chi_s\cong \pi_2 X$, which imply the desired equality $\pi_2 X \cong \varprojlim_s \pi_2 X\cong \varprojlim_s \pi_2 \chi_s=0$.
%\end{proof}

%\begin{theorem}\label{nth100}
%Assume $\overline{W}X$ be a finite type subcontractible presentation assotiated with a free simplicial group presentation $X$ of $\pi_0X$ then for each $\pi_2 \mathbb{Z}_s \overline{W}X\cong\pi_2 \overline{W}X/\Gamma_s X =0$.
%\end{theorem}
%\begin{proof}
%Our finiteness assumption on $X$ imply that  $\mathbb{Z}_s \overline{W}X\simeq \overline{W}X/\Gamma_s X$ is a nilpotent space (the conjugation action of fundamental group is nilpotent and homotopy groups are finitely generated).
%Now it follows from \cite[Theorem ]{HMR} that $\pi_2 \mathbb{Z}_s \overline{W}X$ can be included in the "arithmetic square" of second homotopy groups of $p$-adic and rational completions as follows

%\end{proof}

%\begin{theorem} \label{nth1}
%Assume $\overline{W}X$ be a finite type subcontractible presentation assotiated with a simplicial free group presentation $X$ of $\pi_0X$ then:
%\begin{description}
%\item[1.]{The homotopy fiber $\chi$ of $X\rightarrow \widehat{X}$ is $\#$-contractible, where $\widehat{X}=\{\mathbb{Z}_sX\}_{s\in \mathbb{N}}$ - the Artin-Mazur completion of $X$ with respect to the complete class of finitely-generated nilpotent groups;}
%\item[2.]{$\varprojlim_s^1 H_2(\chi_s)=0$, where $\chi_s$ is a fiber of the fibration $\overline{W}X\rightarrow \overline{W}\mathbb{Z}_sX$;}
%\item[3.]{The intersection $\Gamma=\cap_{i=1}^{\infty}\Gamma_n$ of the lower central series $\Gamma_n$ of $\pi_1\overline{W}X$ coincide with  $\mathcal{P}\pi_1\overline{W}X$ - the maximal perfect normal subgroup of $\pi_1\overline{W}X$;}
%\item[4.]{The natural map $\tau:(\overline{W}X)^+\rightarrow \mathbb{Z}_{\infty}\overline{W}X$ is a covering, corresponding to $\pi_1(\overline{W}X)^+=\pi_1\overline{W}X/\Gamma\leq \pi_1\mathbb{Z}_{\infty}\overline{W}X$. $\mathbb{Z}_{\infty}\overline{W}X$ is homotopy equivalent to a 2-complex and therefore $(\overline{W}X)^+$ is also homotopy equivalent to a 2-complex.}
%\end{description}
%$$\xymatrix @R=0.5cm{
%                &         (\overline{W}X)^+  \ar[dd]^{\tau} \ar[dr]^{\kappa}    &\\
 %\overline{W}X \ar[ur]^{()^+} \ar[dr]_{\mathbb{Z}_{\infty}}       & &     \mathbb{Z}_{\infty}\overline{W}X_{\pi_1\overline{W}X/\Gamma} \ar[ld]^{\widehat{\tau}}    \\
 %               &         \mathbb{Z}_{\infty}\overline{W}X    &}$$
%\end{theorem}
%\begin{proof}
%1. The inverse system $X\rightarrow \widehat{X}$ is represented by the cofinal system of maps $X\rightarrow \mathbb{Z}_sX$ 
%\cite[Proposition 3.3.]{Dror73} 
%which induce epimorphisms of fundamental groups and thus each object in the inverse system of $\chi$ is connected. Denote by $\widetilde{\widehat{X}}$ the universal covering of $\widehat{X}$, and by $\widetilde{X}$ the induced covering of $X$. By \cite[Theorem (4.11)]{AM}
%$\widetilde{\widehat{X}}$ is $\#$-isomorphic with the completion $\widehat{\widetilde{X}}$ of $\widetilde{X}$, and the fibre of $\widetilde{X}\rightarrow \widetilde{\widehat{X}}$ is $\#$-isomorphic with that of $\widetilde{X}\rightarrow \widehat{\widetilde{X}}$. Thus we may suppose $\widetilde{X}$ simply connected, and in fact that the inverse system comprising $\widetilde{X}$ consists of simply connected objects.

%2. We remind that by \cite[Proposition 6.3]{GJ} the functor $\overline{W}:SGr\rightarrow S_0$ from simplicial groups to reduced simplicial sets preserves fibrations. The desired triviality of $\varprojlim_s^1 H_2(\chi_s)=0$ is obvious since $H_2(\chi_s)=0$, being a subcomplex of 2-dimensional subaspherical complex $\overline{W}X$ for which $H_2(\overline{W}X)=0$. 
%Indeed, assume that $\overline{W}X<\overline{W}Y$, where $\overline{W}Y$ is a contractible 2-complex. As $\overline{W}Y$ is contractible it follows that $H_2(\overline{W}Y)\cong =0$ and hence $H_2(\overline{W}X)=0$ since $\overline{W}Y$ is a 2-complex i.e it has no 3-cells.

%3. We derive this assertion from two previous. 
%It follows from 2. that $\varprojlim_s^1 H_2(\chi_s)=0$. Now the Milnor's exact sequence \cite[Proposition 6.14.]{GJ}
%$0\rightarrow \varprojlim^1 H_2(\chi_s)\rightarrow H_1(\varprojlim_s\chi_s) \rightarrow\varprojlim_s H_1(\chi_s) \rightarrow 0$
%in homology imply the isomorphism of abelian groups $\varprojlim_s H_1(\chi_s)\cong H_1(\varprojlim_s\chi_s)$. 
%By 1. $\varprojlim_s H_1(\chi_s)=0$ and hence  $\varprojlim_s H_1(\chi_s)=0$.

%By Theorem \ref{c4} $\pi_2\mathbb{Z}_{\infty}\overline{W}X=0$, now it follows from the long exact sequence in homotopy groups of the homotopy fibration $\overline{W}X\rightarrow \mathbb{Z}_{\infty}\overline{W}X$ with the homotopy fiber $\chi_{\infty}$ that 
%$$\Gamma\cong\pi_1\chi_{\infty}^c\cong Ker(\pi_1\overline{W}X\rightarrow \pi_1\mathbb{Z}_{\infty}\overline{W}X),$$ 
%where $\pi_1\chi_{\infty}^c$ is the fundamental group of the connected component $\chi_{\infty}^c$ of $\chi_{\infty}$. 
%Observe that there is a homotopy equivalennce $\mu:\chi_{\infty}^c\simeq \varprojlim_s\chi_s$. 
%Indeed, consider a natural simplicial map $\mu=\varprojlim_s\mu_s$, where $\mu_s: \chi_{\infty}^c\rightarrow \varprojlim_s\chi_s$. 
%The map $\chi_s$ arises in the following homotopy commutative diagram \cite[???]{BK}
%\begin{equation}\label{ne1}
%\xymatrix @R=0.5cm{
%  \chi_s \ar[dr]& & \\
%   \chi_{\infty}^c \ar[u]^{\mu_s} \ar[r]          &    \overline{W}X \ar[r]^{\mathbb{Z}_{\infty}} \ar[dr]_{\mathbb{Z}_s}   & \mathbb{Z}_{\infty}\overline{W}X\ar[d]^{\tau_s} \\
% &              &      \overline{W}X/\Gamma_s X }
%\end{equation}
%arises from the homotopy equivalences $\mathbb{Z}_{\infty}\overline{W}X\simeq\varprojlim_s \mathbb{Z}_s \overline{W}X\simeq \varprojlim_s\overline{W}X/\Gamma_s X $ \cite[???]{BK}.

%Now the Milnor's exact sequence 
%$0\rightarrow \varprojlim^1 \pi_3(\chi_s)\rightarrow \pi_2(\varprojlim_s\chi_s) \rightarrow\varprojlim_s \pi_2(\chi_s) \rightarrow 0$ for homotopy groups and commutativity of \eqref{ne1}
%imply the isomorphism $\pi_2(\varprojlim_s\chi_s) \cong\varprojlim_s \pi_2(\chi_s)\cong \pi_2 \overline{W}X$. 
%It also follows that $\varprojlim_s^1\pi_k\chi_s=0$ for $k\geq 1$ moreover as a consequence of asphericity we obtain an isomorphism $\varprojlim_s\pi_1\chi_s\simeq \pi_1\chi_{\infty}$ and hence the desired homotopy equivalence.
%As a result $\Gamma/[\Gamma,\Gamma]\cong H_1(\Gamma)\cong H_1(\chi_{\infty}^c)\cong\varprojlim_s H_1(\chi_s)=0 $.
%\end{proof}

\section{Prounipotent presentations and profinite generator spaces}\label{s5}
We use the generator-space and continuity conventions of
\cite[Sections~1--2]{MikhContinuous}.  For a set $S$, let
$S^+=S\sqcup\{*\}$ have isolated points $s\in S$ and cofinite
neighbourhoods of $*$.  A family indexed by $S$ is called convergent
to $1$ when it extends to a continuous pointed map from $S^+$.

In characteristic zero, put
\[
 E_S=k^S,\qquad E_S^\vee=\operatorname{Hom}_{\mathrm{cts}}(E_S,k)
       =k^{(S)}.
\]
Here $E_S$ has the product topology, $k$ is discrete, and $k^{(S)}$
is the direct sum, regarded as a discrete vector space.  Define
\[
 \widehat{\mathbb L}(E_S)
 =\varprojlim_{W,c}
   \mathbb L(E_S/W)/\gamma_{c+1}\mathbb L(E_S/W),
 \qquad F_k(S)=\exp\widehat{\mathbb L}(E_S),
\]
where $W$ runs through open subspaces of finite codimension, $c\geq1$,
and $\mathbb L$ is the ordinary free Lie algebra.  The displayed
quotients are finite-dimensional nilpotent Lie algebras.  On the
$k$-points of a prounipotent group we use the inverse-limit topology
with discrete finite-dimensional unipotent quotients.  The category
itself has prounipotent group schemes as objects and group-scheme
morphisms as maps; normal closures and quotients are taken in this
category.

The canonical generators of $F_k(S)$ are $\exp(e_s)$, where $e_s\in E_S$
is the coordinate vector.  Their images in every finite-dimensional
unipotent quotient are $1$ for all but finitely many $s$.  By the
universal property of the completed free Lie algebra, continuous
pointed assignments $S^+\to U(k)$ extend uniquely to morphisms
$F_k(S)\to U$; see \cite[Lemma~2.2]{MikhContinuous}.

In characteristic $p$ we work with pro-$p$ groups and continuous
homomorphisms.  Write $F_p(S)$ for the free pro-$p$ group on the
convergent family indexed by $S$, characterized by
\[
 \operatorname{Hom}_{\mathrm{cts}}(F_p(S),U)
 \cong\operatorname{Map}_{\mathrm{cts},*}(S^+,U).
\]
We write $F(S)$ for either construction, according to the category;
in formulas common to both cases, $U(k)$ denotes $U$ in the pro-$p$
case.  For infinite $S$, freeness here refers to the specified pointed
profinite generator space $S^+$.  The identifications with degreewise
completions of discrete presentation models used below concern finite
rank in each simplicial degree.

For finite $S$ and characteristic zero, $F_k(S)$ is the inverse limit
of the unipotent completions of $\Phi(S)/\gamma_s\Phi(S)$.
For arbitrary $S$, the coordinate Hopf algebra has the description
\[
 \mathcal O(F_k(S))\cong T(E_S^\vee)=T(k^{(S)}),
\]
with shuffle multiplication and deconcatenation coproduct
\cite[Example~4]{Mikh2023}.  Thus
\[
 \mathcal O(F_k(S))\cong
 \varinjlim_{T\subseteq S,\ |T|<\infty}\mathcal O(F_k(T)),
 \qquad
 F_k(S)\cong\varprojlim_{T\subseteq S,\ |T|<\infty}F_k(T).
\]
The transition maps retain the generators in $T$ and send the others
to $1$.  The corresponding pro-$p$ construction has coordinate algebra
given by locally constant $\mathbb F_p$-valued functions
\cite[Example~6]{Mikh2023}; it likewise satisfies
$F_p(S)\cong\varprojlim_{T\subseteq S,\ |T|<\infty}F_p(T)$.
For the generator spaces, $S^+\cong\varprojlim_T T^+$.
We also use the notation $(\widehat S^*,*)$ for this chosen pointed
profinite space, so that $F(\widehat S^*,*)$ means $F(S)$.
A decomposition into finitely many disjoint subsets gives the
corresponding coproduct in the chosen category; in particular,
$F(X\sqcup Y)\cong F(X)\amalg F(Y)$.

Lifts of an infinite family are always required to be convergent.
The existence of such lifts follows from the projective lifting
statement \cite[Proposition~2.3]{MikhContinuous}: if $q:U\twoheadrightarrow V$
is an epimorphism in the chosen category, every morphism $F(S)\to V$
lifts to a morphism $F(S)\to U$.  In characteristic zero, choose a
continuous linear section of the surjection
$\operatorname{Lie}(U)\to\operatorname{Lie}(V)$ and extend its composite
with the map on $E_S$ by the completed free Lie universal property.
The linear section need not preserve brackets.  In the pro-$p$ case
use projectivity of the free pro-$p$ group.  This produces one
continuous morphism and hence a convergent family of lifts in either
case, with no restriction on $|S|$.

\begin{definition}\label{d7}
Let $F=F(X)$ and let $R$ be a closed normal subgroup in the chosen
category.  Put
\[
 V_R=\begin{cases}
 (R/\overline{[R,F]})(k),&\operatorname{char}(k)=0,\\
 R/\overline{R^p[R,F]},&k=\mathbb F_p.
 \end{cases}
\]
A simplicial presentation with a specified relator basis consists of
a continuous pointed map $\tau:Y^+\to R(k)$ whose image normally
generates $R$, together with a topological linear isomorphism
$k^Y\xrightarrow{\sim}V_R$ sending $e_y$ to the class of
$r_y=\tau(y)$.  Equivalently, these classes are the coordinate
pseudobasis dual to a chosen basis of the discrete space
$V_R^\vee\cong k^{(Y)}$.  With $G=F/R$, the associated diagram is
\begin{equation}
\xymatrix{
 F(X\sqcup Y)\ar@<0ex>[r]^{d_0}\ar@<-2ex>[r]^{d_1}
 &F(X)\ar@<-2ex>[l]_{s_0}\ar[r]^{\pi}&G
}\label{5}
\end{equation}
where
\[
 d_0(x)=d_1(x)=x,\qquad d_0(y)=1,\qquad d_1(y)=r_y,
 \qquad s_0(x)=x.
\]
Continuity of $\tau$ and the universal property of $F(X\sqcup Y)$
make these maps morphisms in the chosen category.
\end{definition}

The basis condition makes $\tau$ injective; since $Y^+$ is compact
and $R(k)$ has the Hausdorff inverse-limit topology, $\tau$ identifies
$Y^+$ with its image as a pointed profinite space.
Continuous lifts of a prescribed coordinate pseudobasis of $V_R$
exist: apply the preceding projective lifting statement to the
quotient $R\to V_R$ and the morphism $F(Y)\to V_R$ defined by the
coordinate vectors.  In characteristic zero, $V_R$ here denotes the
corresponding commutative prounipotent group as well as its $k$-points.
The specified normal-generation and basis conditions are the data
required for the presentation theory of
\cite[Section~4.5]{Mikh2023}.  For the finite discrete subpresentations
used below, the next lemma verifies these conditions for the given
words of relations; continuity of their finite indexing family is
automatic.

\begin{lemma}[Independence of the chosen relators]\label{lem:relator-independence}
Let $(X_1\mid Y_1)$ be a finite subpresentation of a contractible
presentation $(X\mid Y)$.  The images of $r_y$, $y\in Y_1$, span a
free direct summand of $\mathbb Z^{(X_1)}$ with this specified basis.
Let $F=F_k(X_1)$ in characteristic zero, or let $F$ be the free pro-$p$
group on $X_1$ with $k=\mathbb F_p$, and let $R_1$ be the closed normal
subgroup generated by these relators.  There are canonical isomorphisms
\[
 k^{Y_1}\xrightarrow{\ \sim\ }V_1,\qquad
 V_1=\begin{cases}
 (R_1/\overline{[R_1,F]})(k),&\operatorname{char}(k)=0,\\
 R_1/\overline{R_1^p[R_1,F]},&k=\mathbb F_p,
 \end{cases}
\]
which send the coordinate vectors to the classes of the chosen relators.
Moreover, these relators extend to a free topological basis of $F$.
\end{lemma}
\begin{proof}
By Lemma~\ref{l.1.3}, the images of all $r_y$, $y\in Y$, form a basis
of $\mathbb Z^{(X)}$.  Let $L$ be the span of those indexed by $Y_1$.
The coordinate projection $\mathbb Z^{(X)}\to L$ restricts to a
retraction $\mathbb Z^{(X_1)}\to L$, since every $r_y$ with $y\in Y_1$
is a word in $X_1$.  This proves the direct-summand assertion, including
when the ambient contractible presentation is infinite.

Tensoring with $k$ gives an injection $k^{Y_1}\hookrightarrow k^{X_1}$.
Normal generation gives a surjection $k^{Y_1}\twoheadrightarrow V_1$:
conjugation by $F$ is trivial in $V_1$, so the images of the specified
normal generators span it.  In characteristic zero this is a statement
about vector groups; in characteristic $p$ it is a statement about
compact elementary abelian groups.  Since $Y_1$ is finite, its span is
closed in both cases.  The composite
\[
 k^{Y_1}\longrightarrow V_1\longrightarrow\operatorname{Ab}_k(F)
       \cong k^{X_1}
\]
is the preceding injection.  Thus the first map is also injective and
is the claimed canonical isomorphism.

Finally extend these vectors to a basis of $k^{X_1}$ and lift the added
vectors to $F$.  The basis theorem for finite-rank free prounipotent
groups, respectively the pro-$p$ Burnside basis theorem and the Hopfian
property of finitely generated profinite groups, shows that the resulting
family is a free topological basis.  This step takes place entirely
inside the completion of $\Phi(X_1)$; no basis theorem for an infinite
ambient pro-$p$ group is required.
\end{proof}

\begin{remark}
In particular, every finite contractible discrete presentation gives a
simplicial presentation with the specified relator basis of
Definition~\ref{d7}.  Indeed, the relators normally generate
$\Phi(X)$, so their images generate the whole closed normal subgroup
$R=F(X)$ after completion.  Lemma~\ref{l.1.3} gives
$|X|=|Y|=n$ and a unimodular exponent-sum matrix.  Over $k$ this
identifies
\[
 k^Y\xrightarrow{\sim}V_R=\operatorname{Ab}_k(F(X)),
 \qquad e_y\longmapsto\operatorname{ab}_k(r_y),
\]
where $\operatorname{Ab}_k(F)=F/\overline{F^p[F,F]}$ in the
pro-$p$ case.  The family is finite, hence convergent, and the basis
condition holds for the original words of relations.

This condition is the one imposed on the relators in
\cite[Section~4.5, Remark~9]{Mikh2023}.  In that same remark,
\emph{proper} means that inflation $H^1(G,k)\to H^1(F,k)$ is an
isomorphism.  For the full contractible presentation $G=1$, this map
is $0\to k^n$, so the presentation is non-proper whenever $n>0$.
More generally, for a finite subpresentation as in
Lemma~\ref{lem:relator-independence}, its source has dimension
$|X_1|-|Y_1|$ and its target has dimension $|X_1|$; thus it is
non-proper whenever $Y_1$ is nonempty.  The criteria in
\cite[Propositions~15 and~16]{Mikh2023} apply to these non-minimal
presentations with the specified relator basis.
\end{remark}

\begin{remark}\label{r7}
For a finite discrete presentation $(X\mid Y)$, form the rational or
pro-$p$ completion of its free simplicial presentation.  The images
of the given relators are a finite, hence continuous, family of normal
generators.  To apply the presentation theory of \cite{Mikh2023}, one
must additionally verify the specified basis condition of
Definition~\ref{d7}; it does not follow from normal generation alone.
For finite subpresentations of contractible presentations,
Lemma~\ref{lem:relator-independence} supplies the required passage from
the discrete cellular calculation to the completed relator-basis condition.
The free quotient identified below then permits application of the
cohomological asphericity criterion of
\cite[Propositions~15 and~16]{Mikh2023}.
\end{remark}
As in Section \ref{r1.2}, we consider a simplicial prounipotent presentation \eqref{5} as the second step of the Eilenberg-MacLane complex building and denote it by $F^{(1)}_{\bullet}$, we also define the second homotopy group of a prounipotent presentation as $\pi_1(F^{(1)}_{\bullet})$. 
For a connected CW-complex $T$, Kan constructs a free simplicial group
$B_T$ with one non-degenerate generator in degree $n$ for each
$(n+1)$-cell of $T$ and a weak equivalence
$B_T\simeq GS_1T$ \cite[Theorem~6.2]{Kan}.  Thus
\[
 \overline WB_T\simeq \overline WGS_1T\simeq S_1T.
\]
We fix the following model convention for the finite presentations below.
Write $K=K(X_1\mid Y_1)$ for the pointed finite presentation complex and
$P_\bullet=P_\bullet(X_1\mid Y_1)$ for its free simplicial presentation.
Each $P_n$ is a finitely generated free group, and there are no
non-degenerate free generators in degrees greater than one.  The reduced
singular complex $S_1K$ is used only for comparison:
\[
 P_\bullet\simeq G(S_1K),\qquad
 |\overline WP_\bullet|\simeq K.
\]
No degreewise finiteness assertion is made about $S_1K$ or its Kan loop
group.  Finiteness of the CW-complex $K$, finite generation of each $P_n$,
and finite homological type are distinct properties; here the first
ensures the latter two for the chosen presentation model.

For $R=\mathbb Q$ or $\mathbb F_p$, write
$\widehat P_\bullet=(P_\bullet)^{\wedge_R}$ for its dimensionwise
completion.  The dimensionwise completion theorem of Bousfield--Kan
\cite[Proposition~4.1]{BK}, applied to this free presentation model of
finite type, supplies the comparison
\[
 R_\infty K\simeq |\overline W\widehat P_\bullet|.
\]
This use of the completion-model theorem is separate from homotopy
invariance of $R_\infty$: we do not infer that dimensionwise completion
preserves arbitrary weak equivalences of simplicial groups.  Kernels,
Moore complexes and filtrations below are formed on the chosen
presentation model or its explicitly indicated completion.

We record the characteristic-zero facts needed for the finite-dimensional
quotients below.  Our convention for the Zariski closure of a subgroup
of a proalgebraic group is the inverse limit of its Zariski closures in
algebraic quotients, as in \cite[Section~2]{HM2003}.  The equivalence between
unipotent groups and finite-dimensional nilpotent Lie algebras, and its
extension to pro-objects, are recalled in
\cite[Section~3.1]{HainWeight2008}; see also \cite[Section~3]{Hain1993}
and \cite[Propositions~14.32, 14.36 and Theorem~14.37]{MilneAG2017}.

In particular, if $U$ is unipotent and $V\triangleleft U$ is a closed
$k$-subgroup, then $(U/V)(k)=U(k)/V(k)$; this is explicitly recalled in
\cite[Section~7.4, proof of Lemma~7.5]{HM2003} (arXiv version).
For prounipotent groups, the same surjectivity follows by applying
$\exp$ to the surjective quotient map of linearly compact Lie algebras.
Likewise, the image of a morphism of prounipotent groups is closed:
its Lie algebra is the image of a continuous linear map between
linearly compact spaces.  Such an image is closed, since continuous
duality identifies it with the annihilator of the kernel of the dual
map.  For the pro-vector-space conventions, see
\cite[Section~2.1]{HainWeight2008}.

\begin{lemma}[Generation on rational points]\label{lem:unipotent-generation}
Let $U$ be a unipotent algebraic group over a field $k$ of characteristic
zero.  If closed $k$-subgroups $U_i$ generate $U$ as an algebraic group,
then
\[
 U(k)=\langle U_i(k):i\in I\rangle_{\mathrm{abs}}.
\]
\end{lemma}
\begin{proof}
Put $H=\langle U_i(k)\rangle_{\mathrm{abs}}$ and
$\mathfrak u=\operatorname{Lie}U$.  Write $C^r\mathfrak u$ for its lower
central series and $U_r=\exp(C^r\mathfrak u)$.
The images of $\operatorname{Lie}U_i$ span
$\mathfrak u/C^2\mathfrak u$: otherwise their inverse image in
$\mathfrak u$ would be a proper Lie subalgebra, whose exponential would
be a proper algebraic subgroup containing every $U_i$.
Under the exponential correspondence, the image of $U_i(k)$ in
$(U/U_2)(k)$ is precisely the corresponding linear subspace.
Consequently $H\to(U/U_2)(k)$ is surjective.

For every $r\geq2$, $C^r\mathfrak u/C^{r+1}\mathfrak u$ is spanned by
length-$r$ brackets from $\mathfrak u/C^2\mathfrak u$.
Choose lifts in $H$ of their entries.  The leading term of the logarithm
of the corresponding iterated group commutator is the Lie bracket.
Scalar coefficients can be absorbed in one entry, whose lift again
exists by surjectivity on the first quotient.  Products add these
leading terms.  It follows that
\[
 H\cap U_r(k)\longrightarrow (U_r/U_{r+1})(k)
\]
is surjective for every $r$.  Starting with an arbitrary element of
$U(k)$, successively correct it by elements of $H$ along this filtration.
The process ends because $\mathfrak u$ is nilpotent, proving $H=U(k)$.
\end{proof}

\begin{theorem}[Field-completion theorem]\label{main}
Let $(X_1\mid Y_1)$ be a finite simplicial subpresentation of a
contractible simplicial presentation $(X\mid Y)$.  Let $Z$ be a set of
cardinality $|X_1|-|Y_1|$.  For $R=\mathbb Q$ or $R=\mathbb F_p$,
\[
 R_\infty\overline W(X_1\mid Y_1)\simeq K(F_R(Z),1).
\]
Here $F_{\mathbb Q}(Z)$ denotes the rational points of the free
prounipotent group over $\mathbb Q$, and $F_{\mathbb F_p}(Z)$ denotes
the free pro-$p$ group.  In particular, the fundamental group is
isomorphic to $F_R(Z)$ and all higher homotopy groups vanish.
\end{theorem}

\begin{proof}
We first treat a field $k$ of characteristic zero.  By
Lemma~\ref{lem:relator-independence}, the specified relators extend to a
free topological basis $Y_1\sqcup Z$ of $F_k(X_1)$, and their images form
the canonical basis of the relative abelianization of their closed normal
closure.  Consequently,
\[
 F_k(X_1)/(Y_1)_{F_k(X_1)}\cong F_k(Z).
\]
The quotient has Hochschild cohomological dimension at most one (zero
if $Z$ is empty).  The cohomological asphericity criterion
\cite[Propositions~15 and~16]{Mikh2023}, in its version for non-minimal
presentations with the canonical relator basis, therefore gives
$\pi_2\overline W(X_1\mid Y_1)_k=0$.
For $k=\mathbb F_p$, the same lemma, applied inside the finite-rank free
pro-$p$ group on $X_1$, gives the complementary free basis and the same
conclusion in the mod-$p$ setting.

The criterion gives precisely
\[
 \pi_1(\widehat P_\bullet)
   =\pi_2(\overline W\widehat P_\bullet)=0.
\]
It does not imply that $\widehat P_\bullet$ is contractible, nor does it
by itself imply vanishing of its higher homotopy groups or of the first
homotopy group of its abelianization.  We use the spectral argument below
to establish vanishing of $\pi_n(\widehat P_\bullet)$ for $n\geq2$.
The additional input for that argument is the injectivity of the
abelianized differential, verified separately using
Lemma~\ref{lem:relator-independence}.
The dimensionwise completion model of Bousfield and Kan identifies the
classifying space of the completed simplicial group with the corresponding
completion of the original presentation complex.

We now continue the proof using the lower central series spectral sequence.
Use the fixed model $P_\bullet=P_\bullet(X_1\mid Y_1)$ and its
dimensionwise completion $\widehat P_\bullet$, with
$R=\mathbb Q$ or $R=\mathbb F_p$.  Thus
$\widehat P_\bullet$ is a simplicial free prounipotent group in characteristic zero
and a simplicial free pro-$p$ group in characteristic $p$.  Its augmentation
has image
\[
 \pi_0\widehat P_\bullet\cong F_R(Z),
\]
where $Z$ is the complementary basis obtained above.  By construction,
$\widehat P_n$ is generated by degenerate elements for every $n\geq2$.

To make explicit the role of the low-dimensional criterion, put
\[
 N_2\widehat P=\ker d_0^2\cap\ker d_1^2,
 \qquad K_i=\ker d_i^1\quad(i=0,1).
\]
The prounipotent (respectively pro-$p$) Brown--Loday lemma, in the
presentation setting of \cite[Section~4.3]{Mikh2023}, gives
\begin{equation}\label{eq:completed-brown-loday}
 d_2^2(N_2\widehat P)=\overline{[K_0,K_1]}.
\end{equation}
Here the bar denotes topological closure in the pro-$p$ case and closed
subgroup-scheme closure in the prounipotent case.\footnote{For a topologically finitely generated profinite group $G$
and a closed normal subgroup $H$, Nikolov--Segal prove that the abstract
subgroup $[H,G]$ is closed; in particular, $[G,G]$ is closed.
See N.~Nikolov and D.~Segal, \emph{On normal subgroups of compact groups},
Theorem~1.1, \href{https://arxiv.org/abs/1310.3359}{arXiv:1310.3359}.
Finite generation of $H$ is not required.  This statement does not by
itself establish closedness of $[K_0,K_1]$ for two closed normal subgroups;
we retain the closure in \eqref{eq:completed-brown-loday}.}

We verify the completed identity using common quotients of the
$2$-truncated simplicial object $G=\widehat P_{\leq2}$.
For the simplicial profinite inverse-limit framework, see Quillen
\cite{QUILLEN69}.  We spell out the common-quotient construction here
for both the pro-$p$ and the prounipotent settings.
Choose normal subgroups $U_n\triangleleft G_n$, $0\leq n\leq2$, with
finite quotients in the pro-$p$ case, or closed normal subgroup schemes
with finite-dimensional quotients in characteristic zero.  Set
\[
 V_n=\bigcap_{\substack{0\leq m\leq2\\
              \alpha:G_n\to G_m\ \mathrm{structural}}}
              \alpha^{-1}(U_m).
\]
There are only finitely many structural maps in this truncated simplex
category, including identities.  Thus the quotients $G_n/V_n$ are finite,
respectively finite-dimensional.  Closure under composition gives
$\beta(V_n)\subseteq V_m$ for every structural map
$\beta:G_n\to G_m$.  These are therefore quotients of the whole truncated
simplicial object, not independently chosen quotients of its degrees.
Since $V_n\subseteq U_n$, they form a cofinal system
$q_\lambda:G\twoheadrightarrow G_\lambda$ with
$G\cong\varprojlim_\lambda G_\lambda$.

The images of the degeneracies generate $G_{2,\lambda}$ abstractly:
this is immediate for finite groups and follows from
Lemma~\ref{lem:unipotent-generation} for unipotent groups.  Hence the discrete
Brown--Loday lemma \cite[Lemma~5.7]{BL} applies at every stage.
The kernels also lift: for $\bar x\in\ker d_i\subseteq G_{1,\lambda}$,
choose a lift $x\in G_1$ and replace it by
$x(s_0d_ix)^{-1}$.  The resulting lift belongs to $K_i$, so
$q_\lambda(K_i)=K_{i,\lambda}$.  Consequently, for $z\in N_2G$,
\[
 q_\lambda(d_2z)\in d_2N_2G_\lambda
 =[K_{0,\lambda},K_{1,\lambda}]
 =q_\lambda([K_0,K_1]).
\]
Cofinality now gives $d_2N_2G\subseteq\overline{[K_0,K_1]}$.
Conversely, for $x\in K_0$ and $y\in K_1$, the element
\[
 z(x,y)=[s_1x,s_1y(s_0y)^{-1}]
\]
satisfies $d_0z=d_1z=1$ and $d_2z=[x,y]$.
Thus $[K_0,K_1]\subseteq d_2N_2G$.  The latter image is closed by
compactness in the pro-$p$ case and by the closed-image property recalled before
Lemma~\ref{lem:unipotent-generation} in characteristic zero.  This proves \eqref{eq:completed-brown-loday}, without
asserting abstract generation in the original completed group or
interchanging an image with an inverse limit.

The identity identifies the kernel of the boundary map of the associated crossed
module with
\[
 \ker\!\left(K_0/\overline{[K_0,K_1]}\xrightarrow{\overline d_1}\widehat P_0\right)
 =\frac{K_0\cap K_1}{\overline{[K_0,K_1]}}
 =\pi_1\widehat P_\bullet=\pi_2(\overline W\widehat P_\bullet).
\]
The cohomological criterion applied above makes this kernel zero.
Thus Brown--Loday supplies the identification, and the cohomological
criterion supplies the vanishing.  Neither step by itself establishes
vanishing in degrees $q\geq3$ of the classifying space.

For the higher-degree argument, we next pass to the $R$-abelianization
\[
 A_\bullet=\operatorname{Ab}_R(\widehat P_\bullet),
\]
where $\operatorname{Ab}_{\mathbb Q}(\widehat P)=\widehat P/\overline{[\widehat P,\widehat P]}$ and
$\operatorname{Ab}_{\mathbb F_p}(\widehat P)=\widehat P/\overline{\widehat P^p[\widehat P,\widehat P]}$.  Since
$\widehat P_n$ is generated by degeneracies for $n\geq2$, the simplicial
$R$-module $A_\bullet$ satisfies
\[
 A_n=D_nA\qquad(n\geq2).
\]
Consequently its normalized complex is concentrated in degrees zero and
one:
\[
 N_nA=0\qquad(n\geq2).
\]
At this point degeneracy gives only $\pi_nA=0$ for $n\geq2$.
The group $\pi_1A$ must be considered separately: the already established
vanishing of $\pi_1\widehat P$ does not establish its vanishing.
Although our remaining target is $\pi_n\widehat P$ for $n\geq2$, the
argument that makes the positive-degree $E^1$-terms zero also uses
$\pi_1A=0$.  Indeed, a nonlinear Lie functor can produce higher homotopy
from degree-one homotopy of its input.

Here the hypothesis is specifically that $(X_1\mid Y_1)$ is a finite
subpresentation of a contractible presentation.  We do not assume that
$K(X_1\mid Y_1)$ itself is contractible.  The ambient contractible
presentation supplies the cellular basis calculation in
Lemma~\ref{l.1.3}, and its consequence is
Lemma~\ref{lem:relator-independence}.
The differential $N_1A\to N_0A$ is injective because the images of the
relators $Y_1$ span the direct summand established in
Lemma~\ref{lem:relator-independence}.  Explicitly, $N_1A=R^{Y_1}$,
$N_0A=R^{X_1}$, and this differential is the exponent-sum map
$e_y\mapsto\operatorname{ab}_R(r_y)$.  Its injectivity is proved
before any homotopy-vanishing assertion is used.  Hence, by the Dold--Kan correspondence,
\begin{equation}\label{eq:abelianization-acyclic}
 \pi_nA_\bullet=H_n(NA_\bullet)=0\qquad(n>0).
\end{equation}
This proves the abelianization assertion separately from the
cohomological criterion.  For a presentation without the above
injectivity, the same degeneracy condition alone would not justify
vanishing of the positive-degree $E^1$-terms; additional analysis of
the spectral sequence would be necessary.
Thus $A_\bullet$ is weakly equivalent to the constant simplicial
$R$-module $\operatorname{Ab}_R(F_R(Z))$.

Write $F_s\widehat P_\bullet$ for the dimensionwise closed lower central
series in characteristic zero.  In characteristic $p$, use instead the
Zassenhaus filtration
\[
 F_sG=D_sG=\overline{\prod_{ip^j\geq s}\gamma_i(G)^{p^j}},
 \qquad i\geq1,\ j\geq0.
\]
In both cases $F_1\widehat P=\widehat P$, and the filtration is complete
and separated.  Since every $\widehat P_n$ is free, the associated graded
simplicial object is the free graded Lie object on $A_\bullet$:
\[
 \operatorname{gr}_s\widehat P_\bullet
 \cong L_s^R(A_\bullet),\qquad s\geq1,
\]
where $L_s^{\mathbb Q}$ is the degree-$s$ part of the free Lie algebra and
$L_s^{\mathbb F_p}$ is its restricted analogue.  A weak equivalence of
simplicial $R$-modules is a simplicial homotopy equivalence, and the
dimensionwise functor $L_s^R$ preserves simplicial homotopies.  It follows
from \eqref{eq:abelianization-acyclic} that
\[
 \pi_n\operatorname{gr}_s\widehat P_\bullet
 =\pi_nL_s^R(A_\bullet)=0
 \qquad(n>0,\ s\geq1).
\]

For the simplicial pro-$p$ framework and its associated spectral
sequences, we refer to Quillen, \emph{An application of simplicial
profinite groups} \cite{QUILLEN69}; for the rational setting, see
\cite{Qui5}.  The chosen filtration gives the unstable Adams--Quillen
spectral sequence
\[
 E^1_{n,s}=\pi_n\operatorname{gr}_s\widehat P_\bullet.
\]
We justify passage from this vanishing to the completed group directly
on the quotient tower, making the convergence argument explicit.
Put $Q_s=\widehat P_\bullet/F_s\widehat P_\bullet$, so that $Q_1=1$.
The degreewise surjective maps $Q_{s+1}\to Q_s$ are fibrations of
underlying simplicial groups, with kernels
$F_s\widehat P_\bullet/F_{s+1}\widehat P_\bullet$.
In characteristic zero surjectivity on rational points follows from the
unipotent quotient property.  The long exact sequences of homotopy
groups and the vanishing of $\pi_n\operatorname{gr}_s\widehat P$ for
$n>0$ imply inductively
\[
 \pi_nQ_s=0\qquad(n>0,\ s\geq1).
\]
Completeness identifies $\widehat P_\bullet=\varprojlim_sQ_s$.
Since this is a tower of fibrations, its inverse limit computes its
homotopy inverse limit.  For $n\geq1$, the Milnor exact sequence has
outer terms
\[
 \varprojlim\nolimits_s^1\pi_{n+1}Q_s=0,
 \qquad \varprojlim_s\pi_nQ_s=0.
\]
(The degree-one statement is interpreted in the usual group/pointed-set
form.)  It follows that
\[
 \pi_n\widehat P_\bullet=0\qquad(n>0).
\]
This supplies the required higher-degree vanishing, together with the
degree-one vanishing already obtained from the cohomological criterion.
Thus the vanishing of the positive-degree terms of the spectral sequence
passes to the limit; no appeal to finite type alone as a convergence
criterion is needed.
Consequently
\[
 \widehat P_\bullet\simeq\operatorname{Const}(F_R(Z))
 \quad\text{and}\quad
 \overline W\widehat P_\bullet\simeq K(F_R(Z),1).
\]
The equivalence with a constant simplicial group is not a claim that
$\widehat P_\bullet$ is contractible: its group of components is
$F_R(Z)$ and need not be trivial.
The classifying space of the completed simplicial presentation, and hence
$R_\infty\overline W(X_1\mid Y_1)$, is therefore aspherical.
\end{proof}

\begin{proposition}[Keune comparison in finite type]\label{prop:keune-finite-type}
Let $k$ be a field of characteristic zero, and let
$X_\bullet\to G$ and $Y_\bullet\to H$ be free simplicial
prounipotent resolutions of finite type over $k$.  Here finite type means
that in every simplicial degree the free basis is finite, and the chosen
bases are compatible with degeneracies.  Every morphism $\alpha:G\to H$
extends to a simplicial morphism $X_\bullet\to Y_\bullet$.
Any two such extensions are simplicially homotopic over $\alpha$, with
all components of the homotopy morphisms of prounipotent groups.
In particular, two such resolutions of the same group are simplicially
homotopy equivalent over that group.
\end{proposition}
\begin{proof}
The proof of Keune's comparison theorem
\cite[Theorem~2]{Ke} applies without change in finite type.
Indeed, any assignment of a finite free basis to $k$-points of a
prounipotent group extends uniquely to a morphism of prounipotent groups.
The resolving property of the target provides the required lifts of
compatible boundaries.  In each degree only finitely many basis elements
must be lifted, while the values on degenerate basis elements are already
prescribed.  Thus Keune's inductive constructions of comparison maps and
simplicial homotopies take place in the prounipotent category itself.
Applying the comparison in both directions over the identity, and its
homotopy uniqueness to the composites, proves the final assertion.
For a full continuous lifting proof, including infinite bases compatible
with degeneracies, see \cite[Proposition~3.1 and Theorem~4.2]{MikhContinuous}.
The construction there uses morphisms from free objects to lift complete
compatible boundaries and an explicit simplicial cylinder for homotopies.
\end{proof}

\begin{corollary}\label{cor:prounipotent-model-comparison}
In characteristic zero, the equivalence with the constant simplicial group
in Theorem~\ref{main} can be realized by morphisms of simplicial
prounipotent groups and simplicial homotopies whose components are
prounipotent morphisms.
\end{corollary}
\begin{proof}
The completed presentation $\widehat P_\bullet$ has a finite free basis
in every degree, compatible with degeneracies.  The preceding proof
establishes that its augmentation to $F_k(Z)$ is a resolution.
Since $F_k(Z)$ is free of finite rank, $\operatorname{Const}(F_k(Z))$
is also a free simplicial prounipotent resolution of finite type.
Proposition~\ref{prop:keune-finite-type} applies to these two resolutions.
We use the original completed presentation throughout; no additional
free resolution is introduced.
\end{proof}

\section{Subcontractible presentations}\label{s6}


The arithmetic-square argument is applied at finite nilpotent stages.
Only then do we pass to a homotopy inverse limit.  The following
formulation records the required order of operations.
\begin{proposition}[Arithmetic squares along a tower]\label{cx5}
Let $T_s$ be an inverse tower of connected nilpotent finite-type spaces,
and write $D_s$ for the three-corner diagram
\[
 (T_s)_0\longrightarrow\mathbb Q_\infty\!\left(\prod_p(T_s)_p\right)
 \longleftarrow\prod_p(T_s)_p.
\]
Here $(T_s)_0=\mathbb Q_\infty T_s$ and
$(T_s)_p=(\mathbb F_p)_\infty T_s$.  If $T=\operatorname{holim}_sT_s$,
then there is a homotopy pullback square
\begin{equation}\label{as}
\xymatrix{
 T\ar[r]\ar[d] &\operatorname{holim}_s\prod_p(T_s)_p\ar[d]\\
 \operatorname{holim}_s(T_s)_0\ar[r] &
 \operatorname{holim}_s\mathbb Q_\infty(\prod_p(T_s)_p).}
\end{equation}
No identification of the bottom-right corner with
$\mathbb Q_\infty(\operatorname{holim}_s\prod_p(T_s)_p)$ is assumed.
\end{proposition}
\begin{proof}
Put $P_s=\operatorname{holim}_{I_{pb}}D_s$.
Theorem~\ref{space-arithmetic} and
Lemma~\ref{degreewise-pullback} give canonical isomorphisms
$\pi_nT_s\cong\pi_nP_s$ for all $n\geq1$, compatible with every
transition map.  Thus the comparison is an isomorphism of the
homotopy-group towers, not merely an abstract isomorphism at each
stage.

For clarity, this also describes the passage to the inverse limit
on homotopy groups.  After functorial fibrant replacement of the
towers, Milnor's natural exact sequences, for $n\geq2$, are
\[
\begin{split}
0\longrightarrow\varprojlim\nolimits_s^1\pi_{n+1}T_s
&\longrightarrow\pi_n\operatorname{holim}_sT_s
\longrightarrow\varprojlim_s\pi_nT_s\longrightarrow0,\\
0\longrightarrow\varprojlim\nolimits_s^1\pi_{n+1}P_s
&\longrightarrow\pi_n\operatorname{holim}_sP_s
\longrightarrow\varprojlim_s\pi_nP_s\longrightarrow0.
\end{split}
\]
The stagewise natural isomorphisms identify both outer terms and
hence the middle terms.  No vanishing of these derived limits is
needed for this comparison.  The corresponding degree-one sequence
is interpreted as an exact sequence of groups, and components by
the usual nonabelian $\varprojlim^1$ of the fundamental groups.
Equivalently, including all components and base points, homotopy
invariance of homotopy inverse limits applied to the objectwise
weak equivalence $T_s\to P_s$ gives
$\operatorname{holim}_sT_s\simeq\operatorname{holim}_sP_s$.

Finally, Proposition~\ref{hofubini} gives
\[
 T\simeq\operatorname{holim}_s\operatorname{holim}_{I_{pb}}D_s
 \simeq\operatorname{holim}_{I_{pb}}\operatorname{holim}_sD_s,
\]
which is precisely \eqref{as}.  This uses commutation of two homotopy
limits, not commutation of a homotopy limit with rationalization.
\end{proof}

\begin{proposition}[Noncommutative arithmetic square for free simplicial groups of finite type]
\label{general-free-arithmetic}
Let $P_\bullet$ be an arbitrary free discrete simplicial group of
finite type: $P_n$ is free of finite rank for every $n$.  No nilpotency
assumption is imposed on $P_\bullet$ or on $\pi_0P_\bullet$.
Put $X=\overline WP_\bullet$.  All the following constructions are
performed degreewise.  Put
\[
 N_{s,\bullet}=P_\bullet/\gamma_sP_\bullet,\qquad
 \widehat P_{\mathbb Q,\bullet}=\varprojlim_s(N_{s,\bullet})_0,
 \qquad
 \widehat P_{p,\bullet}=\varprojlim_s(N_{s,\bullet})_p^{\wedge}.
\]
These are respectively the rational prounipotent and pro-$p$
completions of the original discrete free simplicial group.  Define
its adelic simplicial group by
\[
 \mathcal A(P)_\bullet=
 \varprojlim_s\varinjlim_{S\subset\pi,\ S\text{ finite}}
 \left(
 \prod_{p\in S}((N_{s,\bullet})_p^{\wedge})_0
 \times\prod_{p\notin S}(N_{s,\bullet})_p^{\wedge}
 \right).
\]
There are natural weak equivalences
\begin{equation}\label{general-group-reconstruction}
 \widehat P_{\mathbb Z,\bullet}:=
 \varprojlim_sN_{s,\bullet}\ \simeq\
 \operatorname{holim}\left(
 \widehat P_{\mathbb Q,\bullet}\longrightarrow\mathcal A(P)_\bullet
 \longleftarrow\prod_p\widehat P_{p,\bullet}\right)
\end{equation}
in simplicial groups, and consequently
\begin{equation}\label{general-space-reconstruction}
 \mathbb Z_\infty X\simeq
 \operatorname{holim}\left(
 \overline W\widehat P_{\mathbb Q,\bullet}
 \longrightarrow\overline W\mathcal A(P)_\bullet
 \longleftarrow\overline W\!\left(\prod_p\widehat P_{p,\bullet}\right)
 \right).
\end{equation}
The homotopy limit in these formulas is over the three-object cospan.
In particular, its three entries are explicit simplicial groups, or
 their classifying spaces, formed from the same $P_\bullet$.
The order $\varprojlim_s\varinjlim_S$ in the adelic entry is retained.
No asphericity assumption on $X$ is required.  If $X\to\mathbb Z_\infty X$
is a weak equivalence, the left side of
\eqref{general-space-reconstruction} can be replaced by $X$.
\end{proposition}
\begin{proof}
Nilpotent groups enter only as intermediate quotients: in each degree
$n$, the group $N_{s,n}=P_n/\gamma_sP_n$ is finitely generated free
nilpotent.  Proposition~\ref{asnfn} supplies its arithmetic square,
and the factorization in Theorem~\ref{asn} makes the corresponding
square of group nerves a homotopy pullback.  For clarity, this passage can be made with the two-sided bar model:
the rational and pro-$p$ product groups act on the adelic group by
left and right multiplication.  The factorization makes this action
transitive in every degree, and the stabilizer of the identity is
$N_{s,n}$.  The inclusion of this stabilizer into the action groupoid
is therefore an equivalence in every degree, naturally in the
simplicial operators.  Its diagonal is a weak equivalence.  The
two-sided bar construction models the homotopy pullback of the three
classifying spaces.  This proves the required simplicial homotopy
pullback assertion; applying the derived loop functor gives the
corresponding assertion for simplicial groups.

Now take the homotopy inverse limit in $s$ and commute the two homotopy
limits by Proposition~\ref{hofubini}.  Every tower involved has
surjective maps in each simplicial degree and hence consists of
fibrations of simplicial groups.  For the adelic tower, lift each
component while keeping the same finite set $S$; this proves its
surjectivity as well.  Thus these homotopy inverse limits are computed
by ordinary inverse limits.  The filtered colimit defining the adelic
entry is also computed degreewise; on group nerves its injective
transition maps are cofibrations.  Finally, ordinary inverse limits
commute with products.  These observations prove
\eqref{general-group-reconstruction}.  Apply the derived classifying-space
functor and the lower-central tower model for integral completion
\cite[Chapter~IV]{BK} to obtain \eqref{general-space-reconstruction}.
\end{proof}

\subsection{The adelic limit and rational pronilpotent completion}
\label{adelic-explicit}
We give a concrete description of the adelic entry and distinguish it
from an interchange of an inverse limit and a filtered colimit.  In this
subsection we index the nilpotent quotients by their class: for a free
group $F$ of finite rank put
\[
 G_s=F/\gamma_{s+1}F,\quad N_{s,p}=(G_s)_p^{\wedge},\quad
 U_{s,p}=(N_{s,p})_0,\quad N_s=\prod_pN_{s,p}.
\]
This shifts the index by one relative to
Proposition~\ref{general-free-arithmetic}.  By Proposition~\ref{asnfn}
we may regard $N_{s,p}$ as a subgroup of $U_{s,p}$ and set
\[
 A_{s,S}=\prod_{p\in S}U_{s,p}\times\prod_{p\notin S}N_{s,p},
 \qquad A_s=\bigcup_{S\text{ finite}}A_{s,S}=(N_s)_0.
\]
Let
\[
 N_p=F_p^{\wedge}=\varprojlim_sN_{s,p},\qquad
 U_p=\varprojlim_sU_{s,p},\qquad N=\prod_pN_p.
\]
All embeddings and projections here are the natural ones.

\begin{proposition}[Explicit description of the adelic limit]
\label{adelic-completion-description}
There is a natural group isomorphism
\begin{equation}\label{adelic-coordinate-description}
 \mathcal A(F):=\varprojlim_s A_s
 \cong\left\{(u_p)_p\in\prod_pU_p:\ 
 \substack{\text{for each }s,\ u_p\bmod s\in N_{s,p}\\
 \text{for all but finitely many }p}\right\}.
\end{equation}
Here $u_p\bmod s$ denotes the projection to $U_{s,p}$.
Let
\[
 \mathcal B(F)=\bigcup_{S\text{ finite}}
 \left(\prod_{p\in S}U_p\times\prod_{p\notin S}N_p\right),
 \qquad K_s=\ker(\mathcal B(F)\longrightarrow A_s).
\]
Then $\mathcal B(F)\to A_s$ is surjective for every $s$, and
\begin{equation}\label{adelic-restricted-completion}
 \mathcal B(F)/K_s\cong A_s,\qquad
 \mathcal A(F)\cong\varprojlim_s\mathcal B(F)/K_s.
\end{equation}
Thus $\mathcal A(F)$ is the separated completion of the restricted
product $\mathcal B(F)$ for this filtration.
Moreover, if $\gamma_jN$ denotes the \emph{abstract} lower central series,
then
\begin{equation}\label{adelic-rational-completion}
 \mathcal A(F)\cong\widehat N_{\mathbb Q}^{\mathrm{nil}}
 :=\varprojlim_s(N/\gamma_{s+1}N)_0.
\end{equation}
These identifications are natural in homomorphisms between finite-rank
free groups, and hence apply degreewise to the free simplicial groups
in Proposition~\ref{general-free-arithmetic}.
\end{proposition}
\begin{proof}
A compatible family $a_s\in A_s$ determines, for each prime $p$, a
compatible family of $p$-coordinates and hence $u_p\in U_p$.
Membership in $A_s$ is exactly the finiteness condition in
\eqref{adelic-coordinate-description}.  Conversely, that condition
allows a tuple $(u_p)_p$ to be projected to $A_s$ on every level.
These two constructions are inverse homomorphisms.

For fixed $S$, inverse limits commute with products, giving
\[
 \varprojlim_sA_{s,S}
 =\prod_{p\in S}U_p\times\prod_{p\notin S}N_p.
\]
To lift $a_s\in A_{s,S}$ to $\mathcal B(F)$, lift its finitely many
coordinates in $U_{s,p}$ to $U_p$, and all remaining coordinates in
$N_{s,p}$ to $N_p$.  These projections are surjective, as their
transition towers consist of surjections.  The lifts retain the same
finite set $S$.  The first isomorphism of
\eqref{adelic-restricted-completion} follows, and then so does the
second.  The intersection of the $K_s$ is trivial by the coordinate
description, so this is a separated completion.

To prove \eqref{adelic-rational-completion}, first observe that $N$ is
a topologically finitely generated profinite group.  If $x_1,\ldots,x_r$
is a basis of $F$, their diagonal images topologically generate $N$:
every finite continuous quotient is a product of finite $p$-groups for
finitely many distinct primes, and the images of these generators
surject onto every factor.  A subdirect subgroup of a product of finite
groups of pairwise coprime orders is the whole product.

By the theorem of Nikolov--Segal, every term of the abstract lower
central series of a finitely generated profinite group is closed
\cite{NS2007}.  In particular, all $\gamma_jN$ and $\gamma_jN_p$ are
closed.  Coordinate projections give one inclusion in
\[
 \gamma_jN=\prod_p\gamma_jN_p.
\]
For the reverse inclusion, $\gamma_jN$ contains $\gamma_jN_p$ supported
in each individual coordinate, and therefore contains all finite-support
products of these subgroups.  They are dense in the displayed product;
closedness gives the reverse inclusion.

For a free pro-$p$ group, quotienting by its closed lower central term
is the pro-$p$ completion of the corresponding free nilpotent group.
This also follows directly from their common universal property for
maps to finite $p$-groups of class at most $s$.  Consequently,
\[
 N/\gamma_{s+1}N
 \cong\prod_p(N_p/\gamma_{s+1}N_p)
 \cong\prod_pN_{s,p}=N_s.
\]
Proposition~\ref{asnfn} now identifies the rationalization of this
quotient with $A_s$.  Taking inverse limits proves
\eqref{adelic-rational-completion}.  All maps used are induced by
projections, inclusions and rationalization, which also proves naturality.
\end{proof}

\begin{proposition}[Failure of the interchange even in rank two]
\label{adelic-noninterchange}
For $F=F(x,y)$, the canonical map
\[
 \iota:\varinjlim_S\varprojlim_sA_{s,S}
       \longrightarrow\varprojlim_s\varinjlim_SA_{s,S}
\]
is injective but not surjective.  Under the preceding identifications,
it is the inclusion $\mathcal B(F)\hookrightarrow\mathcal A(F)$.
\end{proposition}
\begin{proof}
The map sends the class of a compatible family $(a_s)_s$ represented
at a fixed finite $S$ to the same family in $\varprojlim_sA_s$.
Two representatives can be compared at the union of their finite sets
of primes.  Since $A_{s,S}\to A_s$ is injective for every $s$, equality
of their images implies equality of the representatives.  Thus $\iota$
is injective.

Choose distinct primes $p_j$ for $j\geq2$, and for each $j$ choose a
basic Hall commutator $c_j\in F$ of weight $j$.  Define $a_s$ by
\[
 (a_s)_{p_j}=(c_j\bmod\gamma_{s+1}F)^{1/p_j}\in U_{s,p_j},
\]
and put all other prime coordinates equal to $1$.  Roots exist uniquely
in these rational nilpotent groups.  When $j>s$, this coordinate is $1$,
so $a_s\in A_{s,\{p_2,\ldots,p_s\}}$.  Projections preserve unique
roots; hence $(a_s)_s$ is a compatible family in $\mathcal A(F)$.

At level $s=j$, however, $c_j^{1/p_j}$ does not lie in $N_{j,p_j}$.
Indeed, the central layer of weight $j$ in $N_{j,p_j}$ is the free
$\mathbb Z_{p_j}$-module on the basic commutators of weight $j$.
The element in question has coordinate $1/p_j$ at $c_j$.  Its projection
to class $j-1$ is trivial by uniqueness of roots, so if it lay in
$N_{j,p_j}$ it would belong to that central integral layer, which is
impossible.  Every fixed set $S$ representing $(a_s)_s$ would therefore
have to contain all $p_j$.  No finite $S$ does, proving non-surjectivity.
The constant simplicial group on $F(x,y)$ gives the same example within
free simplicial groups of finite rank in every degree.
\end{proof}

\begin{remark}[The adelic model and its naturality]
The two steps in the construction have different roles.  The
augmentation-completion description of Quillen supplies the
nilpotent tower computing group rational pronilpotent completion.
Proposition~\ref{adelic-completion-description} then identifies its
terms and its limit explicitly with the pro-$p$ adelic model.
In particular, it proves the natural isomorphism
$\widehat N_{\mathbb Q}^{\mathrm{nil}}\cong\mathcal A(F)$.

On every nilpotent stage there is also the standard natural
identification
\[
 \mathbb Q_\infty BN_s\simeq B(N_s)_0=BA_s.
\]
Since the maps $A_{s+1}\to A_s$ are surjective, passage to the tower
and the classifying-space construction gives
\[
 \operatorname{holim}_s\mathbb Q_\infty BN_s
 \simeq B\!\left(\varprojlim_sA_s\right)=B\mathcal A(F).
\]
This is the space represented by the adelic entry.  The exceptional
finite set of primes may depend on the nilpotent stage, as expressed
by \eqref{adelic-coordinate-description};
Proposition~\ref{adelic-noninterchange} shows that replacing it by a
single finite set changes the limit.

The arithmetic square uses this explicit tower model.  The canonical map
\[
 \mathbb Q_\infty BN\longrightarrow\operatorname{holim}_s\mathbb Q_\infty BN_s
\]
is a separate space-completion comparison; its being a weak equivalence
is not needed or asserted in the construction above.
\end{remark}

\begin{theorem}[Main theorem]\label{c4}
Let $(X_1\mid Y_1)$ be a finite subpresentation of a contractible
presentation $(X\mid Y)$, and let $|Z|=|X_1|-|Y_1|$.  For
$R\in\{\mathbb Q,\mathbb F_p,\mathbb Z\}$ one has
\[
 R_\infty S_1K(X_1\mid Y_1)\simeq K(F_R(Z),1),
\]
where the field cases use the notation of Theorem~\ref{main}, and
\[
 F_{\mathbb Z}(Z)=\widehat{\Phi(Z)}_{\mathrm{nil}}
 =\varprojlim_{s\geq2}\Phi(Z)/\gamma_s\Phi(Z).
\]
\end{theorem}
\begin{proof}
The field assertions, including the identification of the fundamental
groups, are Theorem~\ref{main}.  For the integral assertion, let
$P_\bullet$ be the free discrete simplicial presentation of
$K(X_1\mid Y_1)$.  Its generators in degrees above one are degeneracies.
Lemma~\ref{l.1.3} shows that the abelianized words of relations form a
primitive independent family in $\mathbb Z^{X_1}$.  Complete this family
to an integral basis and choose words $z\in\Phi(X_1)$ lifting the added
basis vectors, indexed by $Z$.

Fix $s\geq2$.  The substitution sending formal generators $u_y$ to the
words of relations $r_y$, and the complementary generators to the
chosen words $z$, is an isomorphism on abelianizations.  On each
lower-central quotient of a free group it induces the corresponding
isomorphism of free Lie powers.  Induction through the finite central
series therefore shows that it is an isomorphism on the free nilpotent
quotient of class at most $s-1$.

Consequently $P_\bullet/\gamma_sP_\bullet$ can be written as the free
nilpotent simplicial group of class at most $s-1$ on a pointed simplicial
set with constant vertices $Z$ and, for each $y$, an edge from $u_y$ to
the base point.  Its edge satisfies $d_0y=1$, $d_1y=u_y$.
Contract these edges to the base point while fixing $Z$.  Applying the
free nilpotent group functor degreewise gives
\[
 P_\bullet/\gamma_sP_\bullet\simeq
 \operatorname{Const}G_s,\qquad G_s=\Phi(Z)/\gamma_s\Phi(Z).
\]
These equivalences and contractions are compatible with $s$, since the
same words and the same edge contractions were chosen at every stage.
In particular,
$T_s:=\overline W(P_\bullet/\gamma_sP_\bullet)\simeq K(G_s,1)$.

The lower-central tower model for integral Bousfield--Kan completion
\cite[Chapter~IV]{BK} now gives
\[
 \mathbb Z_\infty S_1K(X_1\mid Y_1)
 \simeq\operatorname{holim}_sT_s
 \simeq\operatorname{holim}_sK(G_s,1).
\]
The quotient maps $G_{s+1}\to G_s$ are surjective.  Their nerves form a
tower of Kan fibrations, so its ordinary limit computes its homotopy
limit.  Since the nerve of a group commutes with limits, the last space
is $K(\varprojlim_sG_s,1)$, as required.

We can also describe this tower arithmetically, retaining the order of
limits.  Apply Proposition~\ref{asnfn} to $G_s$ and write
\[
 A_{s,S}=\prod_{p\in S}U_{s,p}(\mathbb Q_p)
          \times\prod_{p\notin S}(G_s)_p^{\wedge},\qquad
 A_s=\varinjlim_S A_{s,S}.
\]
Here $U_{s,p}$ is the unipotent $\mathbb Q_p$-group associated with
$G_s$.  This notation refers to nilpotent groups at a fixed stage;
there is no scalar extension from $\mathbb F_p$ to $\mathbb Q_p$.
Because $G_s$ is a finitely generated torsion-free nilpotent group,
$(G_s)_p^{\wedge}$ embeds in $U_{s,p}(\mathbb Q_p)$.  Thus all transition
maps $A_{s,S}\to A_{s,S'}$, $S\subseteq S'$, are injective, and their
nerves are monomorphisms, hence cofibrations of simplicial sets.
One may use a cofinal sequence of finite sets of primes.  The ordinary
colimit then computes the homotopy colimit.  Since the nerve commutes
with filtered colimits,
\[
 \operatorname{hocolim}_S BA_{s,S}\simeq
 \varinjlim_S BA_{s,S}=B\!\left(\varinjlim_S A_{s,S}\right)=BA_s.
\]
This cofibration assertion is made in simplicial sets, not in simplicial
groups.

Likewise, the transition maps of each of the towers
\[
 G_s,\qquad (G_s)_0,\qquad \prod_p(G_s)_p^{\wedge}
\]
are surjective.  Their nerves are therefore towers of Kan fibrations,
so their homotopy inverse limits are computed by ordinary inverse
limits.  Define the concrete groups
\[
 F_{\mathbb Z}(Z)=\varprojlim_sG_s,\qquad
 F_{\mathbb Q}(Z)=\varprojlim_s(G_s)_0,\qquad
 F_{\mathbb F_p}(Z)=\varprojlim_s(G_s)_p^{\wedge}.
\]
The latter two agree with the free prounipotent and free pro-$p$ groups
used above.  For the pro-$p$ identification, every finite $p$-group
quotient of $\Phi(Z)$ is nilpotent and hence factors through some $G_s$.
The rational identification is the lower-central construction of free
prounipotent completion.  Ordinary limits commute with products, giving
\[
 \varprojlim_s\prod_p(G_s)_p^{\wedge}
 =\prod_pF_{\mathbb F_p}(Z).
\]

The adelic tower also has surjective transition maps.  Indeed, an
 element of $A_s$ belongs to $A_{s,S}$ for some finite $S$.  Each of its
 components lifts at stage $s+1$: the rational unipotent maps and the
 pro-$p$ maps are surjective.  These lifts belong to $A_{s+1,S}$, with
 the same set $S$.  Thus $A_{s+1}\to A_s$ is surjective.  Consequently,
\[
 \operatorname{holim}_s BA_s\simeq B\mathcal A(Z),\qquad
 \mathcal A(Z):=\varprojlim_s\varinjlim_S A_{s,S}.
\]
The order of the two ordinary limits in this definition is essential;
no interchange is asserted.

For this particular presentation, the passage to the limit may also
be read directly from Milnor's sequence.  The compatible contractions
above give $T_s\simeq K(G_s,1)$ at every stage, so, for $n\geq2$,
\[
 0\longrightarrow\varprojlim\nolimits_s^1\pi_{n+1}T_s
 \longrightarrow\pi_n(\operatorname{holim}_sT_s)
 \longrightarrow\varprojlim_s\pi_nT_s\longrightarrow0
\]
has both outer terms zero.  Moreover,
$\pi_1(\operatorname{holim}_sT_s)\cong\varprojlim_sG_s$;
the limit is connected because the tower of $G_s$ is surjective.
The same observations apply to the towers of group nerves in the
other entries.  Here the vanishing follows from the explicit
stagewise models and surjectivity, not from an assertion that every
pro-object has vanishing derived inverse limits.


The factorization assertion in Theorem~\ref{asn} makes the homotopy
pullback at each nilpotent stage connected, with fundamental group
$G_s$ and no higher homotopy.  Proposition~\ref{cx5} therefore gives the
following \emph{noncommutative arithmetic square}:
\begin{equation}\label{as2}
\xymatrix{
 \operatorname{holim}_s T_s\ar[r]\ar[d] &
 B\!\left(\prod_pF_{\mathbb F_p}(Z)\right)\ar[d]\\
 BF_{\mathbb Q}(Z)\ar[r] & B\mathcal A(Z).}
\end{equation}
Here $B$ denotes the nerve of an abstract group.  The upper-left corner
is the homotopy pullback of the other three corners, which are now
classifying spaces of explicit groups (equivalently, the
$\overline W$-constructions on constant simplicial groups).
The maps are induced by the compatible maps at the nilpotent stages.
The preceding argument additionally identifies the upper-left corner
with $BF_{\mathbb Z}(Z)$, but the displayed form retains its role as
the homotopy limit of the arithmetic diagram.
Finally, $\overline W$ is a right Quillen functor and hence its derived
functor preserves homotopy limits.  This is the precise homotopical
form of the adjunction used in passing between simplicial groups and
classifying spaces; right adjointness alone would only give preservation
of ordinary limits.
\end{proof}

\begin{remark}[Scope of the asphericity argument]
The compatible stagewise contractions used here depend on the
hypothesis of a finite subpresentation of a contractible presentation.
They give the aspherical completions and the explicit arithmetic square
of Theorem~\ref{c4}.  This contraction argument is not asserted for
every subcomplex of an arbitrary aspherical complex.  In contrast,
the arithmetic square of Proposition~\ref{general-free-arithmetic}
applies to every free discrete simplicial group of finite type and
requires no asphericity assumption.
\end{remark}

% The material below records an alternative spectral-sequence approach.
% It is excluded from the present version because it is not needed for the
% main theorem and its convergence argument requires additional hypotheses.
\iffalse
\section{Keune comparison theorem and Brown-Loday lemma in the prounipotent settings}\label{Keune}
In this section we assume that the reader is familiar with \cite[\S 1,2]{Ke} and at least briefly has viewed \cite[Ch I,\S 6]{And}.
By an augmented simplicial group $G$ we mean a simplicial group $X$ together with a group $X_{-1}$ and morphism $d_0:X_0\rightarrow X_{-1}$, such that 
$d_0 d_0 = d_0 d_1: X_1\rightarrow X_{-1}$.
We will also need a concept of simplicial kernel $Z_nX, n\geq -1$ of $(d^n_0,d^n_1,\ldots, d^n_n)$ in $X$, which is a subobject of $X_n^{n+2}$ defined as follows 
$$Z_n X=\{(x_0,\ldots, x_{n+1})\in X_n^{n+2} | d_i x_j = d_{j-1} x_i, i<j\}.$$ 
For each $n$, there is a homomorphism 
$$d=d^{(n)}: X_{n+1}\rightarrow Z_nX,$$
defined by $d(x)=(d^n_0(x),d^n_1(x),\ldots, d^n_n(x))$ for $x\in X_{n+1}$.
An augmented simplicial group will be called aspherical if for each $n\geq -1$ the homomorphism $d^{(n)}$ is surjective and we say that $X$ is a resolution of $X_{-1}$. 
If each object $X_n, n\geq 0 $ is a free group (a pro-$p$-group or a prounipotent group), then we will call $X$ a free simplicial resolution if moreover the free generators $A_n\subset X_n$ can be chosen in such a way that $s_iX_n\subset X_{n+1}$ for all $0\leq i\leq n$. 

These concepts are linking to the calculations of homotopy groups and we shall start with the following resuld of Keune \cite{Ina}.
\begin{proposition}\label{np2}
If $X$ be an augmented simplicial group then there is 
$\tilde{\gamma}_ n: \pi_n(X) \rightarrow Z_n(X)/dG_{n+1}$ - a natural isomorphism, arising in the diagram as follows
$$\xymatrix{1\ar[r] &d^{n+1}_{n+1} NX_{n+1} \ar[r]\ar[d]& \widetilde{X}_n \ar[d]_{\gamma_n} \ar[r]&\pi_n(X)\ar[r]\ar[d]_{\tilde{\gamma}_n}& 1 \\
1\ar[r] &dX_{n+1} \ar[r]& Z_nX \ar[r]&Z_nX/dX_{n+1}\ar[r]& 1},$$
where 
$\widetilde{X}_n= \{ x\in X_n: d(x)=1 \}, \gamma (x)= (1, \ldots, x)$ and $NX$ is the Moore complex of $X$.
\end{proposition}

Analogousely with the discrete groups situation \cite[6, Ch. I]{And} and \cite[1]{Ke}, for any  prounipotent group $G$ there is a free prounipotent simplicial resolution of $G$.
Such a resolution can be constructed step by step in the same manner as in the discrete case \cite[6, Ch. I]{And} using a simplicial analog of the glueing cells procedure (along the boundary maps presenting non-zero elements of homotopy groups). 
Moreover there is an analog of "a comparison theorem" \cite[Theorem 2]{Ke} at least for finite type simplicial prounipotent resolutions.

%This resolutions may be constructed step-by-step using the simplicial kernel concept  for finite type prounipotent presentations:
%Step 1. Choose a set of generators $Y\subset X_{-1}=G$ and define $X_0=F(Y_0)$, where $Y_0=\{x^{(0)}|s\in X\}$.
%The surjection $d_0: F(X_0)\rightarrow $

\begin{theorem}[Keune comparison theorem]\label{nt1}
Assume we are given free prounipotent augmented simplicial groups  $X$ and $Y$ with finitely generated (as prounipotent groups) $X$ and aspherical $Y$ and let $\alpha: X_{-1}\rightarrow Y_{-1}$ be a prounipotent group homomorphism.
Then
\begin{description}
\item[(1)]{There exists a simplicial prounipotent groups homomorphism $\gamma: X\rightarrow Y$ such that $\gamma_{-1}=\alpha$ ($\gamma$ is said to extend $\alpha$).}
\item[(2)]{If also $\gamma'$ extends $\alpha$, then there exists a simplicial homotopy $h:\gamma\simeq\gamma'$, where $h_i$ are homomorphisms of prounipotent groups.}
\end{description}
\end{theorem}
\begin{proof} The proof is a compilation of Keune's proof of \cite[Theorem 2]{Ke}, using finiteness assumtion for $X$. 
In the free prounipotent groups case we have the universal property of finitely generated free prounipotent group $X_i$ with respect to maps of generators to prounipotent groups. 
Indeed free prounipotent groups are the prounipotent completions of the corresponding discrete free groups of the same rank. 

$X$ being free, there exists so the called $CW$-basis, i.e. subsets $K_n\subset X_n$ of free generators, and the degeneracy operators $s_i$ take $K_n $ into $K_{n+1}$.
Necessary homomorphisms will be defined by the maps from $K_n$.

Construction of $\gamma$.

$\gamma_{-1}$ : $\gamma_{-1}=\alpha$.

$\gamma_0$  : Let $x\in K_0$. Since $Y$ is aspherical, $d_0: Y_0\rightarrow Y_{-1}$ is surjective. Hence there exists a $y\in Y_0$ such that $d_0 y=\gamma_{-1} d_0 x$. Define $\gamma_0 x=y$.

$\gamma_1$ : Take $x\in K_1 - s_0 K_0$. Then $(\gamma_0 d_0 x, \gamma_0 d_1)\in Z_0 Y$. Since $Y$ is aspherical there is a $y\in Y_1$ with $dy=(\gamma_0 d_0 x, \gamma_0 d_1)$. We define $\gamma_1 x=y$. For degenerated elements $s_0 x\in K_1$, where $x\in K_0$ we define 
$\gamma_1(s_0 z)=s_0 \gamma_0 z$, then the required identities of a simplicial space homomorphism hold $d_0\gamma_1=\gamma_0 d_0$ and $d_1\gamma_1 = \gamma_0 d_1$.

$\gamma_n$ : 

Assume $\gamma_{-1}, \gamma_0, \ldots,\gamma_{n-1}$ are already defined and all relevant identities hold. Let $x\in K_n-\cup^{n-1}_{i=0}s_i K_{n-1}$. 
Then $(\gamma_{n-1} d_0 x, \ldots, \gamma_{n-1} d_n x)\in Z_{n-1}Y$. By asphericity of $Y$ there is  $y\in Y_n$ such that $$dy=(\gamma_{n-1} d_0 x, \ldots, \gamma_{n-1} d_n x)$$. 
Define $\gamma_n(x)=y$. 
For $s_i x\in s_i K_{n-1}$ define $\gamma_n (s_i z)=s_i\gamma_{n-1}z$. It is s straightforward to check that $\gamma_n$ is well defined on $K_n$ a satisfies the simplicial identities.

Construction of simplicial homotopies is the same as in the proof of \cite[Theorem 2]{Ke}. 
\end{proof}

By Definition \cite[4.7]{Mikh2023}, a prounipotent simplicial presentation $X$ is aspherical if
\[
 \pi_2 X=\frac{\ker d^1_0\cap\ker d^1_1}
 {\overline{[\ker d^1_0,\ker d^1_1]}}=0,
\]
where the bar denotes topological closure in the pro-$p$ case and closed
subgroup-scheme closure in the prounipotent case.  The Brown--Loday lemma
below identifies this closed commutator subgroup with the Moore boundary
$d_2(NX_2)$ and therefore shows that the truncated simplicial group
\begin{equation}\label{new3}
\xymatrix{
& F(s_1^1 s_0^0(X)\cup s_0^1(Y)\cup s_1^1(Y)) \ar@<1ex>[r] \ar@<-1ex>[r]_(0.6){d_0^2, d_1^2,d_2^2} \ar[r] & {F(s_0^0(X)\cup Y)} \ar@<-3ex>[l]_(0.4){s_0^1,s_1^1} \ar@<-2ex>[l] \ar@<0ex>[r] \ar@<-1ex>[r]_(0.6){d_0^1,d_1^1} &
F(X) \ar@<-1ex>[l]_(0.4){s_0^0} \ar[r] & G}
\end{equation}
is aspherical as a truncated simplicial group, i.e. $\pi_2 F_{\bullet}=0$.

\begin{lemma}[Brown--Loday]\label{nl2}
Let $G$ be a simplicial pro-$p$ group, or a simplicial prounipotent group
over a field of characteristic zero.  Suppose that $G_2$ is generated, in
the corresponding topological or group-scheme sense, by the degenerate
subgroups $s_0G_1$ and $s_1G_1$.  Put
\[
 N_2G=\ker d_0^2\cap\ker d_1^2,
 \qquad K_i=\ker d_i^1\quad(i=0,1).
\]
Then
\[
 d_2^2(N_2G)=\overline{[K_0,K_1]},
\]
where the right-hand side is the closed subgroup generated by the
commutators $[x,y]$, $x\in K_0$, $y\in K_1$.
\end{lemma}

\begin{proof}
We first verify the inclusion which is also responsible for the choice of
indices.  For $x\in K_0$ and $y\in K_1$, set
\[
 z(x,y)=[s_1x,\,s_1y(s_0y)^{-1}]\in G_2.
\]
The simplicial identities give
\[
 d_0z(x,y)=d_1z(x,y)=1,
 \qquad d_2z(x,y)=[x,y].
\]
Thus $z(x,y)\in N_2G$ and
\(,[K_0,K_1]\subseteq d_2(N_2G),).  The latter image is closed.  Indeed,
in the pro-$p$ case $N_2G$ is compact and its continuous image in the
Hausdorff group $G_1$ is compact, hence closed.  In characteristic zero,
the image of a morphism of (pro)unipotent group schemes is a closed
(pro)unipotent subgroup scheme.  Consequently,
\[
 \overline{[K_0,K_1]}\subseteq d_2(N_2G).
\]

For the reverse inclusion, the argument of \cite[Lemma~5.7]{BL} uses only
the simplicial identities and the fact that $G_2$ is generated by
degenerate elements.  Applied to the underlying abstract subgroup generated
by $s_0G_1$ and $s_1G_1$, it gives
\[
 d_2(N_2G)\subseteq [K_0,K_1]
\]
before passage to closure.  If the degenerate subgroups generate only
topologically (respectively, as a proalgebraic group), the same calculation
gives the inclusion into \(\overline{[K_0,K_1]}\), since both sides are
closed.  Equivalently, one may apply the discrete Brown--Loday calculation
to the compatible finite pro-$p$ quotients, respectively finite-dimensional
unipotent quotients, and then pass to the inverse limit.  This proves the
opposite inclusion and hence the asserted equality.
\end{proof}


\section{Several spectral sequences and their applications}\label{Spectral}
For brevity we will use the notation $F_R(Z)$ for free $R$-prounipotent group when $R=\mathbb{Q}, \mathbb{F}_p$ (a pro-$p$-group is indeed a group scheme over $\mathbb{F}_p$ \cite[4.1]{Mikh2023}) and for free pronilpotent group when $R=\mathbb{Z}$.

When $G$ is a prounipotent affine group scheme over a field of characterietics zero we denote by $Ab(G)$ its quotient by the Zariski closure of the commutant $Ab(G)=G/[G,G]$ and when $G$ is a pro-$p$-group $Ab(G)=G/G^p[G,G]$ - the quotient  by $mod(p)$ commutant, i.e. the closure of $G^p[G,G]$ in the pro-$p$-topology.

We will need the following theorem of D. Quillen \cite{QUILLEN} regading spectral sequences of bisimplicial groups
\begin{theorem}\label{qt} If $G=G_{nm}$ is a double simplicial group, then there are two spectral sequences
$$E^2_{pq}=\pi^h_p\pi^v_q G => \pi_{p+q}(\triangle G)$$
$$E^2_{pq}=\pi^v_p\pi^h_q G =>\pi_{p+q}(\triangle G),$$
where $(\triangle G)_n=G_{nn}, d_j=d_j^hd^v_j, s_j=s^h_j s^v_j$.
\end{theorem}

The cohomology groups $H^*(G,M)$ of a prounipotent group with coefficients in $G$-module $M$ can be computed  explicitly using the \textbf{Hochschild complex} \cite[3.2]{Mikh2023}.
We will need a simplicial topological linerly-compact vector spaces continuous dual to the Hochschild complex (we refer the reader to \cite[3.4]{Mikh2023} regarding "continuos duality"). 
Since $Hom({}_{-},k)$ is an exact functor on the category of vector spaces, the homotopy groups of this dual complex will coincide with $H^*(G,M)^*$. We refer the reader to\cite[3.4]{Mikh2023}
\begin{equation}\label{hoh8}{
C^n(G,M)^*=M^*\widehat{\otimes}^n \mathcal{O}(G)^*, \widehat{\partial}^n:C^{n+1}(G,M)^*\rightarrow C^{n}(G,M)^*,}
\end{equation}
where $d^{n+1}=\sum^{n+1}_{i=0}(-1)^i\widehat{\partial}_i^n$ for  $n\in \mathbb{N}$ is defined by formulas
$$d_0^{n+1}(m\widehat{\otimes} f_1\widehat{\otimes}\ldots\widehat{\otimes} f_n)=\Delta_M(m)^*\widehat{\otimes} f_1\widehat{\otimes}\ldots\widehat{\otimes} f_n,$$
$$d_i^{n+1}(m\widehat{\otimes} f_1\widehat{\otimes}\ldots\widehat{\otimes} f_n)=m\widehat{\otimes} f_1\widehat{\otimes}\ldots\widehat{\otimes} f_{i-1}\widehat{\otimes}\Delta_G(f_i)^* \widehat{\otimes} f_{i+1}\widehat{\otimes} \ldots\widehat{\otimes} f_n,\mbox{ for $1\leq i \leq n$},$$
$$d_{n+1}^{n+1}(m\widehat{\otimes} f_1\widehat{\otimes}\ldots\widehat{\otimes} f_n)=m\widehat{\otimes} f_1\widehat{\otimes}\ldots\widehat{\otimes} f_n\widehat{\otimes} 1,$$
$$s_i^{n+1}(m\widehat{\otimes} f_1\widehat{\otimes}\ldots\widehat{\otimes} f_n)=(m\widehat{\otimes} f_1 \widehat{\otimes} \ldots f_i \widehat{\otimes} 1 \widehat{\otimes} f_{i+1} \ldots \widehat{\otimes} f_n).$$
\begin{proposition}\label{np1}
 Let $X_{\bullet}$ be a finite type free prounipotent simplicial resolution of a prounipotent group $G$, then $\pi_n Ab(X) \cong H^{n+1}(G,k_a)^*$ and $\pi_n Ab(X) \cong H^{n+1}(G,\mathbb{F}_p)^*\cong  H_{n+1}(G,\mathbb{F}_p)$
in the pro-$p$-case.
\end{proposition}
\begin{proof}
Consider the bisimplicial topological linearly-compact vector space $S_{ij}=\widehat{\otimes}^{j-1} \mathcal{O}(X_i)^*$, 
with the horizontal simplicial structure of the free simplicial resolution $X_i$ of $G$.
The vertical structure of $S_{ij}$ comes from the simplicial structure in the dual Hochschild complex of $X_i$.

(1) If we fix $j$ we obtain a horizontal simplicial topological linerly-compact vector space $\widehat{\otimes}^{j-1} \mathcal{O}(X_i)^*$, where $i$ varies. 
Since $X_i$ - is a free simplicial resolution (it is aspherical), hence $X_i^k$ is also aspherical and therefore $\widehat{\otimes}^{j-1} \mathcal{O}(X_i)^*$ is also aspherical, by duality, since for any $k>0$ there is an isomorphism $\mathcal{O}(X_i^k)\cong \otimes^k \mathcal{O}(X_i)$ of Hopf algebras of regular functions of affine group schemes.
%$$\pi_0(\widehat{\otimes}^{m-1} \mathcal{O}(X_n)^* )\cong \widehat{\otimes}^{m-1} \mathcal{O}(G)^*, n=0$$
%$$\pi_n(\widehat{\otimes}^{m-1} \mathcal{O}(X_n)^* ) = 0, \, n \ge 1.$$
%Homotopy groups of $C^m(G,k_a)^*$ are dual to the cohomology groups of $G$ and we obtain
%$$\pi_m^v \pi_0^h (Y_{nm}) = H^m(G,k_a)^*.$$
Now the second Quillen's spectral sequence takes the form
$$E_{n,m}^2 = \pi_n^v \pi_{m}^h (S_{ij}) =\begin{cases}
H^{n+1}(G,k_a)^*&\text{if $m=0$},\\
0&\text{if $m \ge 1$.}
\end{cases}
$$

(2) If we fix $i$ we obtain a vertical simplicial topological linerly-compact vector space $C^n(X_i,k_a)^*$.
It's homotopy groups coincide with $\pi_n(C^n(X_i,k_a)^*)\cong H^n(X_i,k_a)^*$. Since $X_i$ are free, we conclude that only the last nonzero term is $k_a$ in dimension zero column, $Ab_R(X_i)$ in dimension one, it is trivial in all other dimensions.
In the notions of the first Quillen spectral sequence this means 
$$E_{n,m}^2 = \pi_n^h \pi_m^v (S_{ij}) = \begin{cases}
\pi_n Ab(X),&\text{m=1}\\
k_a,&\text{m=0,n=0},\\
0, &\text{if $m \ne 1$ or $(m=0,n\ne 0)$},
\end{cases}
$$
where $\Psi(X)=[X,X]$ if $char(R)=0$ and $\Psi(X)=X^p[X,X]$ if $char(R)=p>0$.

By Quillen's Theorem \ref{qt} we obtain that $\pi_n(X/\Psi(X)) \cong \pi_{n}(\widehat{\otimes}^{j-1} \mathcal{O}(G)^*)\cong H^{n+1}(G)^*$.
\end{proof}
A free restricted Lie algebra on a $\mathbb{F}_p$-module $M$ and ket $TM$ be the tensor algebra $TM=\sum_{r\geq 0} M^{\otimes r}$.
For $a,b\in TM$ we define $[a,b]=ab-ba$ and $a^{[p]}=a^p$. The free restricted $p$-Lie algebra $L^p M$ on $M$ is the smallest sub $\mathbb{F}_p$-module of $TM$ containing $M$ and closed under the operations $[_{-},_{-}]$ and $_{-}{}^{[p]}$.
Denote $L^p_r=L^p\cap M^{\otimes r}$ so that $L^p=\sum_{r\geq 1} L^p_r M$.  For each $r$, $L^p_r$ is a functor of $M$.

If $G$ is a group let $\Gamma_r^p\; G$ be its $p$-lower series and let $gr^p\; G=\oplus \Gamma_r^p G/\Gamma_r^{p+1} G$ be the associated graded $p$-Lie algebra over $\mathbb{F}_p$.
Denote $L^p V= \oplus L^p_r V$ the free Lie algebra generated by a $\mathbb{F}_p$-module $V$, then the canonical map 
$$L^p(gr^p_1\; G)\rightarrow gr^p \; G$$
of $p$-Lie algebras which extends the identity map in degree 1 is always an epimorphism and is an isomorphism for $G$ is a free group by Zassenhaus result.
If $G$ is a free simplicial group, then the component-wise filtration $\Gamma_r^p G$ gives rise to the $p$-lower central series spectral sequence. Indeed  the $p$-lower central series filtration $\Gamma_r^p G$ leads to the homotopy exact couple of the filtered group complex(see \cite{Rec} for basic constructions ragarding the unstable Adams $mod(p)$-spectral sequence) and therefore to the spectral sequence as follows
$$E^1_{nm}=\pi_n L^p_m( gr_1^p\; G)\hspace{10mm} d_r: E^r_{nm} \rightarrow E^r_{n-1, m+r}$$
Suppose we are given a free simplicial group which is finetely generated in each dimension, then we are in the situation of the Quillen's result \cite[Proposition 1.11]{QUILLEN69} below

\begin{proposition}\label{q1}
If $G$ is a free simplicial group such that $H_q(G,\mathbb{F}_p)=\pi_{q-1}(gr^p_1)$ is of finite dimension for each $q$, then the spectral sequence converges strongly to $\pi_q G^{\wedge}_p$ - the homotopy groups of the pro-$p$-completion of $G$.
\end{proposition}
We will need the following rational analoque of the previous proposition also in a situation of finitely generated free simplicail prounipotent groups. The rational Adams spectral sequence apears in \cite[]{Qui5} and is defined analogousely with the $mod(p)$ version as the rationalized Curtis spectral sequence \cite{Cur}
$$E^1_{nm}=\pi_n L_m^{\mathbb{Q}}(gr_1 G\otimes\mathbb{Q}) \hspace{10mm} d_r: E^r_{nm} \rightarrow E^r_{n-1, m+r}$$ 
\begin{proposition}\label{q2}
If $G$ is a free simplicial group such that $H_q(G,\mathbb{Q})=\pi_{q-1}(gr_1\; G\otimes \mathbb{Q})$ is of finite dimension for each $q$, then the spectral sequence converges strongly to $\pi_q G^{\wedge}_{\mathbb{Q}}$ - the homotopy groups of the prounipotent completion of $G$.
\end{proposition}
\begin{proof}
Define the lower central series filtration for a prounipotent group $G$ by the rule $\Gamma_r G=\varprojlim \Gamma \Gamma_r G/G'$, where $G'$ runs over the Zarisski closed normal subroups of finite codimesio, i.e. $G/G$'is an unipotent group.  The category of prounipotent groups is equivalent to the category of pro-finite dimenisonal nilpotent Lie algebras. 
The last are indeed pro-finite dimensional vector spaces and therefore satisfies the Mittag-Leffler condition (we assume that groups are countably generated, since they arise from $CW$-complexes which are constructed by the inductive procedure), so this condition also satisfies prounipotent groups.
And therefore inverse limit functor is exact for prounipotent groups $gr\; G=\varprojlim G/G'$ and $G=\varprojlim G/\Gamma^r G$.
If $V$ is a pro-finite-dimensional vector space over $\mathbb{Q}$, then define its Lie algebra as $L(V)=\varprojlim L(V/V')$, so that we obtain a canonical map $L(gr_1 \; G)\rightarrow gr\; G$.
The last map is an isomorphism for arbitrary sets of generators (not only finite) \cite[Definition 10 and Remark 6]{Mikh2016}.
The map $G\rightarrow G^{\wedge}_{\mathbb{Q}}$ carries $\Gamma_rG$ into the Zariski normal closured lower central normal subgroup $\Gamma_r G^{\wedge}_{\mathbb{Q}}$ and therefore induces the map of spectral sequences 
$$\pi_n(gr_k G\otimes \mathbb{Q})\rightarrow \pi_n (gr_k G^{\wedge}_{\mathbb{Q}}), k\geq 1$$
$gr_1\; G \otimes \mathbb{Q}$ is a simplicial $\mathbb{Q}$-module hence is homotopy equivalent to $B=\sum_g K(\pi_q gr_1\; G\otimes \mathbb{Q})$ - the simplicial $\mathbb{Q}$-module, which is of finite dimension by hypothesis.

By Theorem \ref{} $(gr_1\; G)^{\wedge}_{\mathbb{Q}}\cong gr_1\; G\otimes \mathbb{Q}\cong gr_1 \; G^{\wedge}_{\mathbb{Q}}$ and hence this simplicial $\mathbb{Q}$-module is homotopy equivalent to $B^{\wedge}_{\mathbb{Q}}$. The fact that the prounipotent completion of a free group is a free prounipotent imply the isomorphism of rational Lie algebras $gr\; G^{\wedge}_{\mathbb{Q}}\cong L(gr_1 \;G^{\wedge}_{\mathbb{Q}})$ and therefore we have the homotopy equivalence $gr\; G\cong L(B^{\wedge}_{\mathbb{Q}})$. Thus we have an isomorphism of spectral sequences. As a consequence of exactnes of $\varprojlim$ in the category of prounipotent groups the second spectral sequence converges strongly to $\pi_n G^{\wedge}_{\mathbb{Q}}$ \cite[]{CA}, and the proposition is proven.
\end{proof}

We will need the following Lemma from \cite[4.4]{BOUSFIELD}
\begin{lemma}\label{nl1}
Let $f:Y\rightarrow Y'$ be a Lie map between  $p$-restricted or rational simplicial free Lie algebras. If $\pi_* Ab f$ is an isomorphism then so is $f_{\pi_*}:\Gamma_s Y/\Gamma_{s+1} Y\rightarrow \Gamma_s Y'/\Gamma_{s+1} Y'$
\end{lemma}
\begin{proof} Denote $Lie: Mod_R\rightarrow Lie_R$ the functor which sends $R$-module $M$ into the free rational unrestricted Lie algebra $Lie^{\mathbb{Q}}(M)$ over $\mathbb{Q}$, when $R=\mathbb{Q}$, and sends sends $\mathbb{F}_p$-module $M$ into the free $p$-restricted Lie algebra $Lie^p(M)$, when $R=\mathbb{F}_p$. 
It also sends a morphism $f: M\rightarrow N$ into the morphism $Lie^{\mathbb{Q}}(f):Lie(M)\rightarrow Lie(N)$ \cite[Ch.II, 5]{Bou} and  to the morphism $Lie^p(f):Lie^p(M)\rightarrow Lie^p(N)$ in the $p$-restricted case \cite[Exercise 4, Ch.II, 4]{Bou}. 
Since simplicial $R$-modules constiture a model category it follows that the map, which induces an isomorphism of homoty groups has a homotopy inverse map \cite[]{Qui1}.
For each $n$ let $H_i, i=1,n$ be a homotopy (\cite[Ch. III, Lemma 2.15.]{GJ}) between the maps of $R$-modules $f_0\circ f_1\simeq id_N$ and $f_1\circ f_0\simeq id_M$ then by \cite[Satz 1.15]{DP} (see also \cite[Ch.II, 1.7-1.8, Ch II, \S 6]{Qui1} in the simplicial category context) $Lie(H_i)$ is a homotopy between $Lie(f_0)\circ Lie(f_1)$.
This homotopy induces a homotopy $gr(Lie(H)): gr(Lie(M))\rightarrow gr(Lie(N))$ of graded objects and therefore the isomorphisms of the required homotopy groups.
\end{proof}

\begin{theorem}\label{nt2}
Assume we are given a simplicial subcontractible presentation $X$, then  for $R=\mathbb{F}_p, \mathbb{Q}$ the Bousfield-Kan completions $R_{\infty}\overline{W}X$ and $R_{s}\overline{W}X$ are the Eilenberg-MacLane spaces.
\end{theorem}
\begin{proof}
We will denote for brevity $F_R^{\wedge}(Z)$ the free prounipotent group over the set $Z$, when $R=k$ - the field of characteristic zero and the free pro-$p$-group over $Z$, when $R=\mathbb{F}_p$ - the prime field of characteristic $p>0$, as well as $X^{\wedge}_R$ are the corresponding dimension-wise prounipotent or pro-$p$-completions of simplicial (or truncated simplicial) groups.
It is proved in Theorem \ref{main} that $\pi_1 X^{\wedge}_R=0$ and therefore by Brown-Loday Lemma \ref{nl2} the degenerated (in dimension greater than 1) also truncated in dimensions greater than 2, the free simplicial prounipotent group $X^{\wedge}_R$
\begin{equation}\label{new2}
\xymatrix{
& F_R(s_1^1 s_0^0(X)\cup s_0^1(Y)\cup s_1^1(Y)) \ar@<1ex>[r] \ar@<-1ex>[r]_(0.6){d_0^2, d_1^2,d_2^2} \ar[r] & {F_R(s_0^0(X)\cup Y)} \ar@<-3ex>[l]_(0.4){s_0^1,s_1^1} \ar@<-2ex>[l] \ar@<0ex>[r] \ar@<-1ex>[r]_(0.6){d_0^1,d_1^1} &
F_R(X) \ar@<-1ex>[l]_(0.4){s_0^0} \ar[r] & F_R(Z)}
\end{equation}
is the begining of a free prounipotent (not minimal) simplicial resolution of the free prounipotent group $F_R(Z)$. As it is shown in Theorem \ref{main} the coresponding simplicial prounipotent presentation is proper.
It is also a consequence of Brown-Loday Lemma \ref{nl2}, that if $\pi_2$ of a presentation coincide with $\pi_2 X^{\wedge}_R$, i.e. $\pi_2 X_R=\frac{Ker\; d_0\cap Ker\;d_1}{[Ker\;d_0,Ker\; d_1]}\cong\pi_2 X^{\wedge}_R$ , where  we use the same notation $X^{\wedge}_R$ for the continuation of the truncated simplicial structure in \eqref{new2} to dimensions greater than 2. 
Such continuation exists, it is also finitely generated, and could be given by the formulas of Andre \cite[Ch I, \S 6, p.28]{And}. 
Now we want to show that $X^{\wedge}_R$ is indeed an Eilenberg-MacLane space of type $K(F_R(Z),0)$, i.e. we need to check that $\pi_k X^{\wedge}_R=0, k\geq 1$.

First we note that, since the second prounipotent cohomology group is trivial $H^2(F_R(Z),R)=0$ for free prounipotent group $F_R(Z)$, it follows from Proposition \ref{np1} that $\pi_1 Ab(X^{\wedge}_R)\cong H^2(F_R(Z),R_a)^*$,
and therefore the (mod-$p$) and the "rational" abelianizations $Ab(X^{\wedge}_R)$ have trivial first second homotopy groups. It also follows from the Dold-Kan theorem \cite[Ch. 3, Theorem 2.1]{GJ} that 
$\pi_k Ab(X^{\wedge}_R)=\pi_k(Ab(X^{\wedge}_R)/D(Ab(X^{\wedge}_R)))=0$ for $k\geq 2$, where the degenerated part $D(Ab(X^{\wedge}_R))$ coincide with the whole group $Ab(X^{\wedge}_R)$ in dimensions greater than 2 by construction.

Since $F_R(Z)$ is free, it follows that it also has the constant free prounipotent simplicial resolution of the form $Const(F_R(Z))$ with identity maps in the simplicial structure. 
By Keune's Theorem \ref{nt1} there is $\gamma:X^{\wedge}_R\rightarrow Const(F_R(Z))$ which extends the identity map 
$\gamma_{-1}: (X^{\wedge}_R)_{-1}=F_R(Z)\rightarrow F_R(Z)=Const(F_R(Z))_{-1}$
$$\xymatrix{X^{\wedge}_R\ar[rr]^{\gamma} \ar[d]^{Ab}&&Const(F_R(Z)) \ar[d]^{Ab} \\
Ab(X^{\wedge}_R)\ar[rr]^{\tilde{\gamma}} &&Const(Ab(F_R(Z)),}$$
where $\tilde{\gamma}$ is the induced morphism of the abelianizations.

Observe that $\tilde{\gamma}$ induces an isomorphism in homotopy groups and therefore, since the simplicial prounipotent groups involved are fre. Since the free algebra functor It also induces an isomorphism on graded quotients 
This map induces the map of the corresponding $mod(p)$ (in the pro-$p$-case) and lower central (when $char(R)=0$) graded objects of the filtrations. 
It follows from Lemma \ref{nl1} that the homotopy groups $\pi_k (\Gamma_s X/\Gamma_{s+1} X)\cong \pi_k L_s^R Ab(X^{\wedge}_R)=0$, where $L_s^R$ is the $s$-th graded guotient of the simplicial unrestricted Lie alsgebra over the simplicial vector space $Ab(X^{\wedge}_R)$ in the case of characteristic zero and the $s$-graded guotient of the $p$-restricted Lie algebra, when $R=\mathbb{F}_p$.

By Propositions \ref{q1},\ref{q2} the unstable $R$-Adams spectral sequences $E^1_{nm}=\pi_n L^R_m (gr_1^p X^{\wedge}_R)\cong\pi_n L^R_m Ab(X^{\wedge}_R)$ converge to homotopy groups of $X^{\wedge}_R$. 
By the above $E^1_{nm}=0, m,n>0$ and hence $\pi_nX^{\wedge}_R=0$ for $n>0$, therefore $X^{\wedge}_R\simeq K(F_R(Z),0)$ are the Eilenberg-MacLane spaces.
\end{proof}
The following Corollary is follows directly from Theorem \ref{nt1} and Lemma \ref{nl2}, since endofunctors on the category of groups applied dimensionwise to simplicial groups respect homotopies.

\begin{corollary}\label{nc1}
Assume we are given a finite type subcontractibe simplicial presentation $X$, then 
\begin{description}
\item[(1)]{The prounipotent completion of the free prounipotent resolution $X_{\bullet}$ (which is the expansion of $X$ to higher dimensions) is homotopy equivalent  $(X_{\bullet})^{\wedge}_R\simeq Const(F_R(Z))$ to the constant simplicial prounipotent group $Const\; F_R(Z)$;}
\item[(2)]{There is a homotopy equivalence $(X_{\bullet})^{\wedge}_R/\Gamma_s (X_{\bullet})^{\wedge}_R\simeq Const\; F_R(Z)/\Gamma_s F_R(Z)$ of free simplicial prounipotent groups, where $\Gamma_s$ is the $s$-th lower central series (closed) normal subgroup;}
\item[(3)]{$\pi_1\mathbb{Z}_s X\cong \Phi(Z)/\Gamma_s\Phi(Z)$ - the $s$-nilpotent group and 

$\pi_1\mathbb{Z}_{\infty} X\cong \Phi(Z)^{\wedge}_{nilp}=\varprojlim_s\Phi(Z)/\Gamma_s\Phi(Z) $ - the free pro-nilpotent group of rank $|Z|$.}
\end{description}


\end{corollary}
\begin{proof}
(1) is a content of Theorem \ref{nt2}.

(2) The lower central series quotients endofunctors $\textgoth{G}_s:Gr\rightarrow Gr$ respect homotopies by \cite[Satz 1.15]{DP} and hence by Corollary \ref{nc1} we proceed to the following homotopy commutative diagram
$$\xymatrix @R=0.5cm{
  (X_{\bullet})^{\wedge}_u \ar[rr]^{\simeq}\ar[d]& &Const\; F_u(Z)\ar[d] \\
  \frac{ (X_{\bullet})^{\wedge}_u}{\Gamma_s (X_{\bullet})^{\wedge}_u}\ar[rr]^{\simeq}& &\frac{Const\; F_u(Z)}{\Gamma_sF_u(Z)}}$$
which imply that $\pi_k\frac{ (X_{\bullet})^{\wedge}_u}{\Gamma_s (X_{\bullet})^{\wedge}_u}= 0$, for $k>0$. 

(3) Since $\mathbb{Z}_s X$ is a nilpotent space by (2) we get a pull-back 
For each $s\geq 2$ (since $\mathbb{Z}_s \overline{W} X$ is a nilpotent space) there is a homotopy pull-back diagram
$$\xymatrix{\pi_1\mathbb{Z}_s \overline{W}X \ar[r]\ar[d]  & \prod_{p\in \pi} (\Phi(Z)/\Gamma_s\Phi(Z))^{\wedge}_p\ar[d] \\
 (\Phi(Z)/\Gamma_s\Phi(Z))^{\wedge}_{\mathbb{Q}}\ar[r] &(\prod_{p\in \pi} ((\Phi(Z)/\Gamma_s\Phi(Z))^{\wedge}_p)^{\wedge}_{\mathbb{Q}}}$$
Observe that $\Phi(Z)/\Gamma_s\Phi(Z)$ is isomorphic to the pull-back of that diagram and therefore $\pi_1\mathbb{Z}_s \overline{W}X\cong \Phi(Z)/\Gamma_s\Phi(Z)$.
Since limits commute, we obtain the required isomorphism $\pi_1\mathbb{Z}_{\infty} X\cong \varprojlim_s\Phi(Z)/\Gamma_s\Phi(Z) $.
\end{proof}





% Authors must disclose all relationships or interests that
% could have direct or potential influence or impart bias on
% the work:
%
% \section*{Conflict of interest}
%
% The authors declare that they have no conflict of interest.


% BibTeX users please use one of
%\bibliographystyle{spbasic}      % basic style, author-year citations
%\bibliographystyle{spmpsci}      % mathematics and physical sciences
%\bibliographystyle{spphys}       % APS-like style for physics
%\bibliography{}   % name your BibTeX data base

% Non-BibTeX users please use

\fi

\section*{Conclusion}
For every finite subpresentation of a contractible presentation, the rational,
mod-$p$, and integral Bousfield--Kan completions of the associated presentation
complex are Eilenberg--MacLane spaces of type $K(\pi,1)$.  The argument uses
the canonical independence of the relators, the spectral argument with its
explicit quotient-tower convergence, and the arithmetic square;
the integral argument uses compatible nilpotent-stage models and their
homotopy inverse limit.

\bibliography{Mikhovich_Whitehead_Arxiv1}{}
\bibliographystyle{plain}




\end{document}
